% GENERATED FILE — DO NOT EDIT; authoritative prose is Markdown.
% source: manuscript/orthemma-ortheme-systems-revised-draft.md
% source-sha256: 35ed98f779bf6e6881b24ab100f70336d43e55ca6ebdd1a5128895b08d285a02
% source-qualification: research-stage-draft
% source-qualification: not-peer-reviewed
% source-qualification: not-externally-peer-reviewed
% source-qualification: not-empirically-validated
% source-qualification: candidate-terminology-not-adopted
\documentclass[10pt,letterpaper,twocolumn]{article}
\usepackage{amsmath}
\usepackage{amssymb}
\usepackage{booktabs}
\usepackage{geometry}
\usepackage{hyperref}
\usepackage{microtype}
\usepackage{natbib}
\usepackage{xcolor}
\geometry{letterpaper,margin=0.75in}
\hypersetup{colorlinks=true,linkcolor=blue,urlcolor=blue,citecolor=blue}
\setlength{\emergencystretch}{3em}
\title{Orthemma–Ortheme Systems: An Analysis-Relative Architecture for Auditable Classification and Handling}
\author{}
\date{}
\begin{document}
\twocolumn[
\maketitle

\begin{quote}
\textbf{Provenance.} Multi-model draft lineage with successive review passes; original headers preserved verbatim in \href{\detokenize{../docs/provenance/document-history.md}}{\texttt{docs/\allowbreak{}provenance/\allowbreak{}document-\allowbreak{}history.\allowbreak{}md}}; nothing here is experimentally validated.

\end{quote}
\textbf{Revised manuscript draft — July 2026 (revision R6; fresh-session repository review completed; not external human peer review; not empirically validated). Not peer reviewed.}


\section{Abstract}
This paper proposes and analyzes an orthemma–ortheme system: an information-processing architecture that encounters concrete occurrences, infers which repeatable operational state-types they instantiate, preserves uncertainty when that identity is unresolved, routes each occurrence by its sufficiently validated profile, and revises its types and governing rules when recurring error shows the current mapping inadequate. The two poles are a token and a type: an \emph{orthemma} is a concrete situated occurrence — this utterance, this word-token, this patient presentation, this build report for this commit — and an \emph{ortheme} is a repeatable operational state-type it instantiates, relative to a declared, versioned \emph{analysis} (of which the task is one component; ground truth is analysis-relative, never task-only and never actor-relative). The thesis is stated at the outset as an \textbf{integration discipline}: every facet of the architecture has an established neighbor (multi-label classification, POMDP belief states, the reject option, value-of-information stopping, risk registers, ticket state machines, provenance records, process reliabilism), and the claimed contribution is not any single facet but the union, the placement lifecycle that connects the facets, and one further object — the reified handling \emph{episode} carrying a joint verdict vector that separates result correctness from pathway adequacy and robustness on a single auditable record. The formal core is extended by six additions: versioned identity with labeled lineage; analysis/actor/time indexing; typed, scoped, expiring evidence channels; factorized candidate structure with route composition; per-burden closure status; and episode reification. The verdict remains conditional: nothing here is experimentally validated. The design was motivated by private internal engineering records that are not independently auditable or published; they function as non-evidential design provenance, and no public claim rests on them. No public observational dataset accompanies the paper. The coined vocabulary must beat matched established terminology on an adequately powered benchmark — externally registered before any run — before any term is adopted.

\textbf{Keywords:} classification lifecycle; type–token distinction; analysis-relative ground truth; auditable episodes; process reliabilism; selective prediction; provenance; pathway adequacy; metamorphic testing; governance.


\bigskip\hrule\bigskip
\vspace{1em}
]

\section{1. Introduction}

\subsection{1.1 Four concrete occurrences}
Consider four concrete occurrences.


\begin{itemize}
\item A person says something, here and now, that reaches a listener as an ambiguous acoustic stream — compatible with "olive juice" and with "I love you."

\item A learner meets an Arabic word-token whose segmentation, lemma, morphology, attached clitics, and syntactic role are not yet settled.

\item A patient arrives vomiting red material; the presentation is compatible with several risk states and several causes.

\item A build system emits "all tests pass" for this commit, which may mean the behaviour was genuinely exercised, or merely that files exist, or that the evidence is stale.

\end{itemize}
Each is a concrete situated occurrence. In none of them is the system's first task to decide whether two \emph{other} cases may be merged. Its first task is to apprehend \emph{this} occurrence: to determine which repeatable, consequence-bearing state-types it instantiates, to act on the part of that identity it has established, and to hold the rest open. That is the phenomenon this paper is about.


\subsection{1.2 Two poles}
We give the phenomenon a type–token vocabulary:


\begin{itemize}
\item \textbf{orthemma} = the concrete situated occurrence;

\item \textbf{ortheme} = the repeatable operational state-type it instantiates.

\end{itemize}
An orthemma–ortheme system is an architecture built around the relation between them: concrete occurrences are encountered, apprehended, investigated, classified, represented, routed, validated, dispositioned, and learned from by determining which orthemes they instantiate. The primary object is the relation and its lifecycle, not an isolated type and not a bare partition boundary.


\subsection{1.3 The thesis, stated as an integration discipline}
\textbf{Central thesis.} An orthemma–ortheme system is an \emph{integration discipline} for the token side of classification and handling work: an information-processing architecture that encounters concrete occurrences, infers which repeatable operational state-types they instantiate, preserves uncertainty when that identity is unresolved, routes each occurrence according to its sufficiently validated orthemic profile, dispositions every residual obligation explicitly, records each handling episode as an auditable object distinct from its result, and revises its orthemes and governing rules when recurring error shows the current mapping inadequate.

We concede up front that each facet of this architecture has an established neighbor. Plural profiles are multi-label and hierarchical classification. The maintained candidate set is a POMDP belief state. Abstention and escalation are the reject option and selective prediction. Deciding when to stop investigating is value-of-information stopping. Residual obligations are risk registers; placement statuses are ticket state machines; lineage is provenance; correct-result-through-defective-process is process reliabilism; perturbation probes are metamorphic testing. None of these facets is claimed as new. What is proposed is threefold: (i) the \textbf{union} — no one of those frameworks, in its standard individual form, supplies the others; (ii) the \textbf{lifecycle} — the explicit connection of disclosure, evidence, representation, routing, validation, closure, and revision to the placement status of a single concrete occurrence and its successors; and (iii) the \textbf{joint episode-verdict object} — a reified handling episode on which result correctness, evidential support, truth-connection, rule adequacy, executor fidelity, route safety, closure truthfulness, and robustness are simultaneously definable and separately adjudicable. All three are an integration proposal with an accompanying utility hypothesis (Section 13), not an established result.


\subsection{1.4 The redundancy critique, answered inline}
The strongest objection is the redundancy critique: \emph{a composite of POMDP belief states, the reject option, multi-label classification, value-of-information stopping, a risk register, and a ticket state machine would act identically; the residual is a glossary.} Our answer, developed in Section 12 and maintained throughout, is: the critique is correct about each facet taken separately, and it is correct that the type-individuation mathematics is inherited. What remains is narrower than impossibility, and still substantive. The six-framework composite the critique invokes is a union that no member framework names: in their standard individual forms, the neighbors do not by themselves supply one shared, occurrence-centered episode/verdict record on which result correctness and pathway adequacy are jointly statable. A custom composite certainly \emph{could} be extended with such a record — this paper's proposal just \textbf{is} that extension, stated once and generically: a common lifecycle, a shared machine-checkable record contract, and the hypothesis that standardizing the integration improves auditability and practice. That hypothesis is adjudicated by the unrun benchmark of Section 13, not by this paper's argument, and no observational claim about how costly the union is to assemble in practice is made here — the synthetic worked example of Section 11.4 illustrates the claimed integration burden without being evidence for it.


\subsection{1.5 Contributions and claim levels}
Relative to the prior manuscript, this revision makes six formal additions:


\begin{center}
\begin{tabular}{@{}p{0.313\linewidth}p{0.313\linewidth}p{0.313\linewidth}@{}}
\toprule
\textbf{\#} & \textbf{Formal addition} & \textbf{What it closes} \\
\midrule
1 & Versioned identity and labeled lineage on the occurrence domain & "right finding, wrong copy"; verdicts silently transported across mutations \\
2 & Analysis / actor / time indices on profiles and candidate sets & ground truth only relative to a declared analysis; cross-actor divergence as a stop signal \\
3 & Typed, scoped, expiring evidence channels & green-but-mis-scoped validators; stale evidence treated as current \\
4 & Factorized candidate structure with route composition & one occurrence, several independent defects, each with its own route \\
5 & Per-burden closure status with residual dispositions & false closure as a type error; scalar "percent complete" banned \\
6 & Episode reification with a joint verdict vector & correct-by-luck placements; pathway adequacy separable from result correctness \\
\bottomrule
\end{tabular}
\end{center}
Five kinds of claim stay separate throughout — conceptual framework, formal definition, empirical hypothesis, metaphysical thesis, and design vocabulary. We defend the first two (the framework's coherence and the formal definitions); the design vocabulary is \emph{proposed but not defended} — every coined term is benchmark-gated (Section 14) and the paper's content survives its removal; the empirical hypotheses are stated with their tests and none has been run (Section 13); and no metaphysical thesis is defended here (the modal and metaphysical material is split to companion papers, with only the ordinary-words relational property — analysis-relative resolvability under perturbation — retained where it motivates the robustness fixtures).


\bigskip\hrule\bigskip

\section{2. Core Objects}

\subsection{2.1 Occurrence, type, and the instantiation relation}
\textbf{Definition 1 (Orthemma).} An orthemma is a concrete situated occurrence — an event, object, utterance, token, episode, or artifact — considered as something to be apprehended: this utterance now, this word-token on the page, this patient episode, this build report for this commit.

\textbf{Definition 2 (Ortheme).} An ortheme is a repeatable operational state-type that an orthemma may instantiate, relative to a declared analysis (Definition 3 and Section 2.6 — the domain and task are among its components). It is a type whose confusion with another type would change a warranted classification, prediction, investigation, route, validation condition, closure condition, or evaluation.

\textbf{Definition 3 (Analysis-relative instantiation).} Let $M$ be a domain of concrete orthemmata (plain gloss: the concrete cases) and $O$ a repertoire of repeatable orthemes (the state-types in play). The primitive instantiation relation is analysis-relative: $\operatorname{Inst}_A \subseteq M \times O$, where $(m, o) \in \operatorname{Inst}_A$ reads: the orthemma $m$ instantiates the ortheme $o$, relative to the declared analysis $A$ — the explicit, versionable index defined in Section 2.6, whose components include the task $T = \operatorname{task}(A)$. (Notation note: the prior manuscript used a task-indexed primitive, now retired — see the notation registry's \texttt{retired\_symbols}; the letter $I$ is reserved for the individuation component of the governed-component key, Section 6. The task-subscripted $\operatorname{Inst}_T$ survives only under the abbreviation convention below.) The \textbf{orthemic profile} of $m$ is its fibre,


\begin{verbatim}
O*(m; A) = { o ∈ O : (m, o) ∈ Inst_A }.
\end{verbatim}
The occurrence $m$ and its worldly facts are not created by the analysis. What is analysis-relative is which profile of consequence-bearing operational state-types truly describes $m$ in the declared repertoire; the actual profile $O^*(m; A)$ remains distinct from the observation of $m$, from the system's inferred profile, from the evidence available, and from any actor's belief.

\textbf{Abbreviation convention (task-relative shorthand).} $O^*_T(m)$ is permitted only as local scoped shorthand: after the text has explicitly fixed one analysis $A$ with $\operatorname{task}(A) = T$, and within that scope only, $O^*_T(m) := O^*(m; A)$ — likewise $\operatorname{Inst}_T$, $\Pi_T$, $\mathcal{K}_T$, $\mathcal{R}_T$, and $\mathcal{W}_T$ for their $A$-indexed counterparts ($\Pi_A$, $\mathcal{K}_A$, $\mathcal{R}_A$, $\mathcal{W}_A$), and the unsubscripted $M$, $O$ for the analysis-active domain $\mathcal{M}_A \subseteq \mathcal{M}$ and repertoire $\mathcal{O}_A \subseteq \mathcal{O}$. The shorthand must not be presented as a globally well-defined function of $T$ alone: two analyses can share one task while differing in tolerance, representation, boundaries, or merger family. It is \textbf{forbidden wherever more than one analysis is live} — higher-order audits, multi-actor evaluation (Section 10), cross-version comparisons, differing tolerances or governance boundaries, and comparisons between base execution and reviewer analysis — where the full $O^*(m; A)$ form is required. \textbf{Standing scope for this paper:} except where a passage explicitly introduces a second analysis, one declared analysis $A$ is fixed with $T = \operatorname{task}(A)$, and task-subscripted notation below is that licensed shorthand, not a second primitive. No separate task-to-analysis bridging law is introduced: there is one ground-truth primitive, and the task-relative form is abbreviation only.

The profile is generally not a singleton. One orthemma instantiates several orthemes at once and at different descriptive levels; the relation may be hierarchical, compositional, overlapping, and temporarily unresolved. A single Arabic word-token can simultaneously instantiate a lexeme, a lemma, a morphological pattern, several affixes or clitics, a syntactic role, a semantic contribution, and a discourse function. A build report can simultaneously instantiate a test-outcome type, an evidence-scope type, and a provenance-currency type. An orthemma is not assigned to one mutually exclusive box.


\subsection{2.2 Observation is not the orthemma}
The occurrence and the signal it presents are different objects. Let $\Omega : M \rightharpoonup X$ be the (partial, typed) observation map and $x = \Omega(m)$ the presently available observation of $m$. The observation may be impoverished, aliased, or misleading; the orthemma is the concrete event that produced or occasioned it. Apprehension runs


\begin{verbatim}
m --Ω--> x --evidence H_t--> p̂_t(m) --placement--> a_t --creates--> Succ ⊆ M
                                                      |______________________|
                                                    (feedback: successors re-enter M)
\end{verbatim}
where $\hat{p}_t(m)$ is the system's current inferred profile (belief), $a_t$ is the resulting interpretation, investigation, route, or action, and — new in this revision — the action's \textbf{successor set} $\operatorname{Succ}$ closes the loop back into the occurrence domain (Section 2.5). The ground truth $O^*_T(m)$ and the belief $\hat{p}_t(m)$ are distinct; so are the observation $x$ and the occurrence $m$.


\subsection{2.3 Known orthemma, unresolved identity}
A system may know that an orthemma occurred while not knowing its orthemic identity: $m$ is encountered, yet $\hat{p}_t(m)$ remains open. The right description is


\begin{verbatim}
m is encountered; p̂_t(m) remains unresolved,
\end{verbatim}
not "an unknown ortheme floats free." Types do not hover independently of occurrences; what is unresolved is the placement of a present, concrete case. This also separates the four things "the hidden distinction" can mean, as facets of one relation: the actual profile $O^*_T(m)$ (ground truth), the evidence and candidate structure (Section 5), the inferred profile $\hat{p}_t(m)$ (belief), and the way the belief is encoded (representation). None entails the others.


\subsection{2.4 Placement and apprehension}
\textbf{Definition 4 (Placement and apprehension).} Placement is the assignment of an inferred profile $\hat{p}_t(m)$ to an orthemma. Apprehension is the whole process of Section 2.2: encountering $m$, observing $x$, accumulating typed evidence, narrowing the candidate structure, placing $\hat{p}_t(m)$, routing $a_t$, and dispositioning what remains. Placement is provisional until validated and may be revised.


\subsection{2.5 Versioned identity and lineage (formal addition 1)}
The prior manuscript said "this build report \emph{for this commit}" and warned that "a correct classification attached to stale provenance is not a valid placement of the current artifact" — but gave those phrases no formal object. This subsection promotes them.

\textbf{Definition 5 (Identity key and version).} Every orthemma carries an identity key $\kappa = \operatorname{id}(m)$ (plain gloss: \emph{which thing this is}, surviving change) and a version $v = \operatorname{ver}(m)$ (\emph{which state or edition of that thing}). Two orthemmata may share $\kappa$ and differ in $v$ (the same file before and after an edit); two may share an observation and differ in $\kappa$ (a reused storage slot occupied by a new file).

\textbf{Definition 6 (Labeled successor edges).} An action $a$ performed in the handling of $m$ creates a labeled successor set


\begin{verbatim}
Succ_a(m) ⊆ M,   each element reached by an edge labeled with the action that produced it,
\end{verbatim}
of size \textbf{zero, one, or many}. (A read-only classification creates none; an edit creates one; a split, broadcast, or build creates many. An earlier formulation forcing a single successor $m' = \operatorname{succ}(m, a)$ is corrected.)

\textbf{Transport principle.} Placement validity is bound to $(\kappa, v)$. A placement or validation established for $(\kappa, v)$ does \textbf{not} transport across a successor edge to $(\kappa, v')$ — or to a different $\kappa$ behind the same observation — without a lineage argument and, where the edge could have changed the placed property, fresh evidence. This makes "right finding, wrong copy" a well-formed, detectable error rather than an invisible one: the finding was valid, and it is attached to an occurrence that no longer exists in the relevant version. Illustrative failure modes of this class — of the kind that motivated the design, stated here generically — include a reused storage slot treated as the same file, verdicts silently carried across a tree substitution, ordinal position mistaken for identity, and a stale checkout evaluated as the current one.


\subsection{2.6 Analysis, actor, and time indices (formal addition 2)}
Ground truth is only defined relative to a \textbf{declared analysis} $A$: the system and governance boundary, task, evidence and action repertoire, policy class, loss, hard constraints, horizon, tolerance, representation family, and permitted-merger family. This section is the definition site of the index Definition 3 takes as primitive: $A$ must be \textbf{explicit and versionable} — it carries an identifier and version $\operatorname{ver}(A)$, and a change to any component above yields a new analysis version. Profiles and beliefs are therefore indexed:


\begin{verbatim}
O*(m; A)       — actual profile under analysis A (ground truth relative to A);
p̂_{A,α,t}(m)   — the profile actor α infers at time t under A;
C_{A,α,t}(m)   — the candidate structure actor α maintains at time t.
\end{verbatim}
There is no second, task-only ground truth alongside this one: $O^*_T(m)$ is Definition 3's scoped abbreviation, licensed only while a single fixed $A$ with $\operatorname{task}(A) = T$ is in force.

Two clarifications guard the index. \textbf{Analysis relativity is not actor relativism:} the analysis is a \emph{declared, versioned, public} index, not anyone's belief state — once $A$ is fixed, $O^*(m; A)$ is an objective matter that every actor can be wrong about, and disagreement between actors is disagreement about one analysis-relative fact, not the coexistence of private truths. \textbf{Analysis-version transport:} a change to any component of $A$ yields $\operatorname{ver}(A)+1$, and placements, validations, and verdicts established under one analysis version do \textbf{not} transport to another without a declared transport argument (which components changed, and why the placed claims are invariant under that change) — the exact discipline Section 2.5 imposes on occurrence versions, applied to the analysis index itself.

Placement status is \textbf{actor-indexed}: a clinician can hold a distinction absent from the record; a developer can know a failure the release gate cannot express; a reviewer can carry a verdict the pipeline has no field for. \textbf{Cross-actor divergence} — two actors' inferred profiles for the same $(\kappa, v)$ disagreeing consequentially — is promoted in this revision from an aside to a first-class stop signal: it is itself diagnostic and is a legitimate trigger for the governed interruption of Section 7.5, independent of either actor's confidence. Multi-actor evaluation is developed further in Section 10.


\subsection{2.7 The orthemic contrast, and the inherited results}
\textbf{Definition 7 (Orthemic contrast).} Fix a declared analysis $A$. Two orthemes are operationally non-equivalent, $o_i ≢_A o_j$, when confusing their concrete instances would alter a warranted classification, prediction, investigation, route, validation condition, closure condition, or evaluation beyond the accepted tolerance. An orthemic contrast is such a non-equivalence: a criterion by which two candidate types must not be collapsed.

The contrast is a relation on types. It explains why two orthemes must remain distinguishable; it does not, by itself, tell a system what a present orthemma is. Contrast individuates the types; instantiation constitutes the system.

Formally, for candidate types realised by evidence histories $h_i, h_j$ with mixture weights $\lambda, 1-\lambda$, let $\operatorname{Rep}_A^{i=j}$ be the family of representations that place the two histories identically (plain gloss: all the ways of treating them as one), and $\operatorname{Rep}_A$ the unrestricted family. The \textbf{merger gap} is


\begin{verbatim}
Δ_A(o_i, o_j) = inf_{χ ∈ Rep_A^{i=j}} L_A*(χ)  −  inf_{χ ∈ Rep_A} L_A*(χ),
\end{verbatim}
with $L_A^*(\chi)$ the best attainable risk under representation $\chi$ and hard-constraint violations counted as infinite risk. Then $o_i ≢_A o_j$ at tolerance $\epsilon_A$ iff $\Delta_A(o_i, o_j) > \epsilon_A$. The individuation is relative to the representation and merger families — components of $A$, which is why the contrast subscript follows the analysis (under Definition 3's convention, $≢_T$ may abbreviate $≢_A$ once a single $A$ is fixed). Another architecture can make the same type distinction feasible or unnecessary.

Three inherited results, kept and presented as inherited (none is a new theorem):


\begin{enumerate}
\item \textbf{Static operational separation.} Force two placements to share an action, and if their best individual actions differ enough, the forced merger incurs excess risk above tolerance — infinite if no common admissible action exists.

\item \textbf{Predictive separation under a strictly proper score.} If the task is forecasting and two placements imply different conditional distributions, one shared forecast is strictly worse than separate truthful forecasts; predictive difference individuates types only when prediction is the declared task.

\item \textbf{Finite-horizon safe merger (controlled bisimulation).} States with identical admissible actions, immediate losses, class-transition probabilities, and terminal losses have equal optimal value; the quotient preserves value. This licenses \emph{just-in-time placement}: leave $\hat{p}_t(m)$ coarse now, run a discriminating test, and refine before the consequential action — preservation of a distinction can be procedural, not a persistent label.

\end{enumerate}
Threshold closeness is non-transitive ($h_1 \approx_{\epsilon_A} h_2$, $h_2 \approx_{\epsilon_A} h_3$, yet $h_1 ≉_{\epsilon_A} h_3$), so there is generally no unique smallest repertoire; several incomparable repertoires can be equally adequate.

\textbf{Definition 8 (Route-sufficient apprehension, retained).} $\hat{p}_t(m)$ is route-sufficient when it pins down enough of $O^*(m; A)$ to select an admissible near-optimal route, even though other orthemes in the profile remain unresolved. It is distinct from identity-complete apprehension, in which the full profile is established for the relevant purpose. Route-sufficiency ≠ complete orthemic resolution. (Hematemesis can alter urgency and routing while the bleeding source stays open; an Arabic token's root and lemma can license dictionary lookup before every clitic is settled.)

\textbf{Definition 9 (Under- and over-segmentation, retained).} A repertoire under-segments an orthemma when it collapses orthemes in $O^*(m; A)$ whose separation the task requires (excess risk above tolerance or a constraint violation). It over-segments when it assigns distinctions the orthemma does not support, or splits at no out-of-sample benefit net of complexity, latency, coordination, fairness, and contestability costs.

\emph{Remark (demotion).} The prior draft's rate–distortion analogue for segmentation is demoted to this remark: its terms were undefined as stated, and a compression framing in any case supplies no semantics and no legitimate ends. It may be reinstated if formalized.


\bigskip\hrule\bigskip

\section{3. Aliasing and Identity Uncertainty}
The type–token framing exposes three related but distinct uncertainty patterns. The prior manuscript distinguished the first two; this revision adds the third as its own column, because it is a distinct pattern in kind — conflatable with, but reducible to, neither of the others.


\begin{center}
\begin{tabular}{@{}p{0.235\linewidth}p{0.235\linewidth}p{0.235\linewidth}p{0.235\linewidth}@{}}
\toprule
\textbf{} & \textbf{Inter-orthemma aliasing} & \textbf{Intra-orthemma uncertainty} & \textbf{Identity uncertainty} \\
\midrule
\textbf{Pattern} & Distinct occurrences produce the same or confusable observation while instantiating operationally different types & One encountered occurrence remains compatible with several candidate types & It is unresolved \emph{which occurrence} (which $\kappa$, or which version $v$) is actually in hand \\
\textbf{Formally} & $\Omega(m_1) = \Omega(m_2)$, $O^*(m_1; A) \neq O^*(m_2; A)$ & $C^{\mathrm{profile}} \supseteq \{p_1, \dots, p_k\}$, several profiles licensing non-equivalent routes & $C^{\mathrm{id}}(m) \subseteq M$ non-singleton: candidates for the case's own identity/version \\
\textbf{Example} & Two jars of white powder look identical; one is flour, one is cornstarch & One acoustic stream compatible with "olive juice" and "I love you" & The log line may describe this commit's build or the previous one's; the file at this path may be the original or a reused slot \\
\textbf{Wrong reduction} & Treating it as one case with an open profile & Treating it as two cases & Treating it as profile uncertainty about a known case \\
\bottomrule
\end{tabular}
\end{center}
The first concerns observational equivalence across concrete cases; the second concerns unresolved placement of a present case; the third concerns unresolved \emph{reference} — the case's own identity key or version is among the open questions. Identity uncertainty is the pattern behind "right finding, wrong copy": there is nothing wrong with the placement as a placement; what is open is which $(\kappa, v)$ it is a placement \emph{of}. A system that maintains only profile uncertainty will misfile identity uncertainty as confidence about the wrong object — which is precisely how verdicts transport across mutations undetected. The typed candidate families of Section 5 give each pattern its own slot.


\bigskip\hrule\bigskip

\section{4. Evidence: Typed, Scoped, Expiring Channels (formal addition 3)}
The prior manuscript's single observation map $\Omega$ is replaced by a family of \textbf{typed evidence channels} $\{\Omega_k\}$, each partial (not every channel applies to every occurrence) and each carrying declared metadata. An evidence item obtained through channel $k$ is a record


\begin{verbatim}
h = ⟨ channel k;  property class τ;  scope σ;  provenance;  validity/expiry ⟩
\end{verbatim}
with components glossed:


\begin{itemize}
\item \textbf{Property class} $\tau$ — the three \textbf{core cross-domain classes} are $\{\mathrm{structural}, \mathrm{behavioral}, \mathrm{provenance}\}$. Structural evidence attests to form (the file parses, the schema matches, the artifact exists); behavioral evidence attests to exercised behaviour (the test ran this code path on this input and observed this output); provenance evidence attests to origin and currency (this artifact was produced by that process from that version). An analysis may declare \textbf{domain-specific subclasses} of these (e.g., histological vs serological within a clinical behavioral/structural scheme); what it may not do is admit evidence carrying \emph{no} declared class, or let one class's pass silently discharge another class's obligation. Exhaustiveness of the three core classes across all domains is a working hypothesis of the framework, not a theorem; the subclass mechanism is the sanctioned extension point. \textbf{Authorization is not an evidence property class.} Whether an actor was \emph{permitted} to place, route, or close is a separate \textbf{warrant gate} $W$ (Section 6), with warrant states such as \{authorized, factually established, both, neither\}: authorization can be present while nothing is established, and a claim can be established that nobody authorized acting on. Treating an authorization record as if it were evidence for the placed claim launders permission into support; the two are never merged.

\item \textbf{Scope} $\sigma$ — the set of claims, occurrence versions, and levels the item can bear on at all (plain gloss: what this check is even \emph{about}).

\item \textbf{Provenance} — which process produced the item, from which $(\kappa, v)$.

\item \textbf{Validity/expiry} — the conditions under which the item remains current; evidence expires when the occurrence it attests to acquires a successor along an edge that could have changed the attested property.

\end{itemize}
\textbf{Green-but-mis-scoped.} A validator can pass while supporting nothing:


\begin{verbatim}
mis-scoped pass  ≡  pass ∧ (σ(h) ∩ claim = ∅).
\end{verbatim}
The check is green, and its scope does not intersect the claim being closed. Observed instances: a hash validator that establishes integrity being read as establishing semantic validity; a stage validator whose scope was the previous stage's output being read as covering the current stage. A mis-scoped pass is not weak evidence — it is \emph{no} evidence for that claim, and under Section 7 it cannot contribute to closure. Symmetrically, stale evidence (valid in its day, expired by lineage) is not weak evidence for the current version; it is evidence about a different occurrence.

Validation discipline follows from the typing: a structural check does not validate a behavioral claim; symptom relief does not establish a cure; a fluent gloss does not establish a correct source-linked parse. Each placed claim declares the property class and scope of evidence that could support it, and an evidence-status map records, per claim: validated, provisional, stale, or absent.


\bigskip\hrule\bigskip

\section{5. Candidate Structure (formal addition 4)}

\subsection{5.1 Typed candidate families}
The prior manuscript's single candidate set $C_t(m)$ is split by uncertainty axis into \textbf{typed candidate families}, indexed like everything else by analysis, actor, time, and descriptive level:


\begin{center}
\begin{tabular}{@{}p{0.313\linewidth}p{0.313\linewidth}p{0.313\linewidth}@{}}
\toprule
\textbf{Family} & \textbf{Ranges over} & \textbf{Open question} \\
\midrule
$C^{\mathrm{id}}$ & $M$ & \emph{which occurrence} (identity/version) is in hand — Section 3's third column \\
$C^{\mathrm{profile}}$ & $\Pi_A$ (profile space) & \emph{which profile} the case instantiates — competing hypotheses may themselves be whole profiles, never coerced into single orthemes \\
$C^{\mathrm{cause}}$ & $\mathcal{K}_A$ (cause repertoire) & \emph{which cause} produced the state \\
$C^{\mathrm{route}}$ & $\mathcal{R}_A$ (route repertoire) & \emph{which operation/owner} should receive the case \\
$C^{\mathrm{warrant}}$ & $\mathcal{W}_A$ (warrant states) & \emph{which warrant state} obtains — authorized, established, both, neither \\
\bottomrule
\end{tabular}
\end{center}
Each family carries an \textbf{exclusivity marking}: elements are flagged as \emph{alternatives} (at most one obtains — this powder is flour or cornstarch, not both) or as \emph{co-holding components} (several may obtain together — this one occurrence has an identity defect \emph{and} a quantity defect). Collapsing co-holding components into alternatives forces a false choice; the reverse collapse hides a disjunction. Optional weights over alternatives are permitted but never required; a family is consequentially open when it contains more than one element whose members would license non-equivalent routes. An open family must carry its \textbf{evidence-to-resolve clause}: which channel, in scope, would discriminate the alternatives. (A list of possibilities with no path to adjudication is treated as a definitional defect here — the norm this definition adopts.)


\subsection{5.2 The profile space and partial profiles}
\textbf{Definition 10 (Profile space, partial profiles — general form first, R3).} Fix a declared analysis $A$ with active repertoire $\mathcal{O}_A$. In its most general form the \textbf{profile space} is


\begin{quote}
$\Pi_A \subseteq \mathcal{P}(\mathcal{O}_A)$ — the set of \textbf{admissible complete profiles} under the constraints declared by $A$,

\end{quote}
where a complete profile is a subset of the repertoire that the analysis's declared constraints admit as a candidate way the occurrence could totally be, at the analysis's level of description. This general form makes no claim that every domain is naturally factorized: hierarchical, compositional, and holistically-constrained profile spaces are all instances (the R3 formal audit and counterexample ledger record the attacks that forced this generality).

\textbf{Factorized representation (one permitted family, not a universal ontology).} An analysis \emph{may} declare its repertoire organized into \textbf{axes} (symptom, cause, severity, evidence-scope, …), each carrying the exclusivity marking of Section 5.1 and an \textbf{applicability condition}. In a factorized representation a complete profile: (i) on each \textbf{applicable} \emph{alternatives}-marked axis selects \textbf{exactly one} admissible value; (ii) on an axis that is \textbf{objectively inapplicable} to the case records the explicit value \texttt{not-applicable} — a declared null, which is a \emph{fact about the case under $A$}, never a representation of uncertainty; (iii) on a \emph{co-holding} axis contains any declared-admissible subset of values, \textbf{including the empty set} where the analysis declares objective absence admissible; and (iv) violates no declared cross-axis \textbf{consistency constraint}. Whether axes partition the repertoire or overlap is itself a declaration of $A$; overlapping axes are admissible only with declared reconciliation constraints. Five states that a factorized representation must never conflate: \textbf{objective absence} (the axis applies; nothing on it obtains — an empty co-holding set), \textbf{objective inapplicability} (the axis does not apply — the \texttt{not-applicable} value), \textbf{epistemic openness} (the actor has not resolved the axis — a property of $\hat{p}$, never of $O^*$), \textbf{evidence absence} (no evidence bears on the axis — a property of $H$), and \textbf{candidate plurality} (several complete profiles remain live — a property of $\hat{C}$). The first two live in profiles; the last three never do.

A \textbf{partial profile} leaves one or more axes (or, in the general form, one or more admissibility questions) undetermined: formally, an assignment of a \emph{set} of still-admissible resolutions (the whole domain when nothing is known), and $\Pi_A^\partial$ is the space of partial profiles; every complete profile is the special case with everything determined, so $\Pi_A \subseteq \Pi_A^\partial$. An \textbf{empty complete profile} (nothing in the repertoire obtains) is admissible exactly where $A$'s constraints admit it — analyses whose repertoires exhaust the possibilities may exclude it by constraint; the framework does not exclude it by fiat.

Four objects now stand in definite relations, and none may be conflated with another:


\begin{itemize}
\item the \textbf{true profile} $O^*(m; A) \in \Pi_A$ — one complete profile, fixed by the occurrence and the analysis, independent of anyone's evidence;

\item the \textbf{candidate set} $\hat{C}_{A,\alpha,t}(m) \subseteq \Pi_A$ (the typed family $C^{\mathrm{profile}}$) — the set of complete profiles the actor has not yet ruled out; correctness of maintenance means $O^*(m; A) \in \hat{C}$ whenever the evidence so far is veridical, and investigation shrinks $\hat{C}$;

\item the \textbf{inferred partial profile} $\hat{p}_{A,\alpha,t}(m) \in \Pi_A^\partial$ — what the actor currently \emph{places}: determined on the resolved axes, open on the rest. A partial profile corresponds to the set of its completions, so $\hat{p}$ and $\hat{C}$ are inter-constrained ($\hat{C} \subseteq \operatorname{completions}(\hat{p})$ when both are maintained honestly) — but \textbf{a candidate set is not one inferred profile}: $\hat{C} = \{p_1, p_2\}$ with two live complete profiles is a different epistemic state from a single vaguer $\hat{p}$, and collapsing the former into the latter loses exactly the alternatives structure that routes and discriminating tests act on;

\item optional \textbf{belief weights} — a distribution over $\hat{C}$ (or over an axis's alternatives). Weights are \emph{permitted, never required} (Section 5.1), and an unweighted candidate set is a legitimate terminal representation, not an unfinished one.

\end{itemize}
This subsection supplies the definition the corpus previously used implicitly ($\Pi_A$ appeared in the candidate-family table without a definition site); it is A-indexed through the repertoire, exclusivity markings, and constraint declarations, all components of $A$ (Decision 0007).


\subsection{5.3 Factorized profiles and route composition}
Profiles factorize over axes: a placement may be resolved on the symptom axis, open on the cause axis, and resolved on the severity axis. Routes then compose. Let $r_1 ⊕ r_2$ denote the \textbf{route composition} of the operations licensed by independently resolved factors, with the semantics:


\begin{itemize}
\item factors with disjoint operational footprints compose freely (treat the identity defect \emph{and} correct the quantity defect);

\item where composed routes conflict, declared priority applies, and \textbf{hard constraints dominate}: a safety constraint on one factor overrides a near-optimal route on another;

\item an unresolved factor contributes no route, but may contribute a containment obligation (Section 7) that composes like a route.

\end{itemize}
This is what makes route-sufficient action well-formed in the plural case: the system acts on the resolved factors' composed route while the open factors persist as explicitly dispositioned residuals rather than being forgotten or blocking all action.


\bigskip\hrule\bigskip

\section{6. Metaorthemes: The Split Normal Form}

\subsection{6.1 Definition}
The prior draft defined a metaortheme as a higher-order distinction "acting on" the system's machinery. This revision adopts a single, narrower normal form, aligned with the core formalization.

\textbf{Definition 11 (Metaortheme, split normal form).} A metaortheme is a \textbf{metaorthemic configuration}


\begin{verbatim}
μ = ⟨ g;  S_μ;  select_μ;  prov(μ);  ver(μ) ⟩
\end{verbatim}
paired with a separable \textbf{meta-policy} $\pi_\mu$, where:


\begin{itemize}
\item $g$ is the \textbf{governed component} — which part of the mapping/handling machinery this distinction governs (Section 6.2);

\item $S_\mu$ is the set of \textbf{declared competing higher-order states} the governing context may occupy, specified \emph{in advance}, not read off one incident (e.g., \{appearance-grade, provenance-grade\} for evidence; \{current, stale\} for version);

\item $\operatorname{select}_\mu$ is the \textbf{selecting evidence} procedure that determines which state in $S_\mu$ actually obtains;

\item $\operatorname{prov}(\mu) = \langle \mathrm{authority}, \mathrm{warrant}, \mathrm{scope}, \operatorname{ver}(\mu)\rangle$ is the rule's own provenance — a rule of unknown provenance is a stale-evidence problem one level up;

\item $\operatorname{ver}(\mu)$ is the rule's version, recorded per episode so audits can scope which placements ran under which edition;

\item $\pi_\mu$, the meta-policy, is the \textbf{conduct rule} that consults the configuration and prescribes behaviour conditional on the obtaining state ("quarantine on stale," "never place on appearance alone").

\end{itemize}
The distinction the rule consults and the rule that consults it are two objects. Two different meta-policies (quarantine vs re-derive) can consult the \emph{same} configuration (current-vs-stale); that shared consultable distinction is what a rules-only account cannot express. A metaortheme \emph{change} induces a transformation of the mapping subsystem — the transformation is the effect of revising $\mu$, not $\mu$ itself.

\textbf{Anti-vacuity conditions (part of the definition, not advice).} A candidate $\mu$ is admitted only if: (i) $g$ is named; (ii) $S_\mu$ is declared in advance; (iii) switching the obtaining state changes validated placement, risk, or constraints in at least one episode class; and (iv) it is not a first-order ortheme, nor a tunable parameter of an ordinary policy in disguise. A good policy that consults no in-advance-declared competing higher-order states is just a good policy. (Whether the \emph{word} "metaortheme" earns keep over ordinary words is a separate question, decided only by the terminology benchmark — Section 14; a real distinction may be admitted while its word is retired.)


\subsection{6.2 Governed components, glossed}
$g \in \{\mathrm{O}, \mathrm{I}, \mathrm{E}, \mathrm{D}, \mathrm{R}, \mathrm{V}, \mathrm{W}\}$, in plain language:


\begin{itemize}
\item \textbf{O — repertoire:} which state-types exist and may be posited or retired;

\item \textbf{I — individuation:} what counts as the same case; identity keys, versioning, lineage;

\item \textbf{E — evidence:} what evidence counts, of which property class, at what grade;

\item \textbf{D — disclosure:} what may be left open, and how open questions must be declared (uncertainty handling);

\item \textbf{R — routing:} where cases are sent, and what happens when the correct route is unavailable;

\item \textbf{V — validation/closure:} what may be called done, at what standard, and when a completion claim must be reopened;

\item \textbf{W — warrant-classification:} which warrant states exist (authorized vs factually established and their combinations) and what each licenses.

\end{itemize}
\textbf{Excluded as governed components:} objectives and loss functions; the task $T$ itself (changing the task changes the problem — rules are task-\emph{indexed}, never task-\emph{governing}); and the governance meta-level (the precedence order among rules and the rights to revise them are parameters fixed at the declared governance boundary, where the regress terminates and which another inquiry may reopen).

\textbf{Negative example (calibration).} In a widely discussed exchange, a conversational system's instruction of the form "be maximally truth-seeking" was described as if it were a governing distinction of this kind. It fails every anti-vacuity condition: it names no governed component, declares no competing higher-order states in advance, specifies no selecting evidence, and is not separable into a consultable distinction and a conduct rule. It is an \emph{objective} — and objectives are excluded from $g$ by construction. The exchange is retained as the standing negative example of an objective mislabeled a metaortheme.


\subsection{6.3 Worked configurations}
Five configurations, each passing the minimum-specificity test:


% breakable-row-table: data-rows=5 total-rendered-characters=744 max-row-rendered-characters=149
\par\medskip
\hrule height 0.8pt
\smallskip
% breakable-row: 1/5
\noindent\textbf{Configuration}: Evidence grade\par
\noindent\textbf{$g$}: E\par
\noindent\textbf{$S_\mu$}: \{appearance-grade, provenance-grade\}\par
\noindent\textbf{$\operatorname{select}_\mu$}: source record or discriminating test\par
\noindent\textbf{Meta-policy examples}: choose which test to run before placing\par
\smallskip
\hrule height 0.4pt
\smallskip
% breakable-row: 2/5
\noindent\textbf{Configuration}: Version currency\par
\noindent\textbf{$g$}: I\par
\noindent\textbf{$S_\mu$}: \{current, stale\}\par
\noindent\textbf{$\operatorname{select}_\mu$}: lineage check of $\operatorname{ver}(m)$ against latest\par
\noindent\textbf{Meta-policy examples}: quarantine on stale; or re-derive on stale\par
\smallskip
\hrule height 0.4pt
\smallskip
% breakable-row: 3/5
\noindent\textbf{Configuration}: Depth of resolution\par
\noindent\textbf{$g$}: R, V\par
\noindent\textbf{$S_\mu$}: \{route-sufficient, identity-complete\}\par
\noindent\textbf{$\operatorname{select}_\mu$}: task declaration + risk class\par
\noindent\textbf{Meta-policy examples}: release the route while holding residuals open\par
\smallskip
\hrule height 0.4pt
\smallskip
% breakable-row: 4/5
\noindent\textbf{Configuration}: Closure standard\par
\noindent\textbf{$g$}: V\par
\noindent\textbf{$S_\mu$}: per-burden dispositions vs "all done"\par
\noindent\textbf{$\operatorname{select}_\mu$}: the residual ledger (Section 7)\par
\noindent\textbf{Meta-policy examples}: claim completion only at the ledger's level\par
\smallskip
\hrule height 0.4pt
\smallskip
% breakable-row: 5/5
\noindent\textbf{Configuration}: Warrant state\par
\noindent\textbf{$g$}: W\par
\noindent\textbf{$S_\mu$}: \{authorized, established, both, neither\}\par
\noindent\textbf{$\operatorname{select}_\mu$}: authorization record vs validating evidence\par
\noindent\textbf{Meta-policy examples}: tag warrant type; never launder one into the other\par
\smallskip
\hrule height 0.8pt
\par\medskip

\subsection{6.4 Plurality, conflict, precedence}
A typical episode runs under several metaorthemes at once — a version rule, an evidence-grade rule, and a closure rule — written $\vec{\mu} = (\{\mu_1, \dots, \mu_k\}, \preceq)$ with $\preceq$ a declared strict partial precedence order. Configurations governing disjoint components compose freely (their constraints conjoin). Two rules conflict at an episode when their meta-policies prescribe incompatible constraints on the same component under the states that actually obtain; resolution is only via $\preceq$, and a conflict unresolved by $\preceq$ is a stop condition (Section 7.5) — silent override is itself a metaorthemic error. Metaorthemes can become first-order orthemes for a higher audit (Section 8.1); the regress stops at the declared governance boundary.


\bigskip\hrule\bigskip

\section{7. Closure, the Lifecycle, and Revision (formal addition 5)}

\subsection{7.1 The lifecycle, recast per burden}
The prior draft drew the lifecycle as one line — encountered → unresolved → candidate profile → investigated → provisionally placed → validated → routed → integrated. The line survives as a \emph{typical trajectory of a single burden}, but the object that matures is not the case as a whole; it is each \textbf{burden} — each claim to be established and each obligation to be discharged — separately. A case has no scalar completion state, and scalar "percent complete" reporting is banned as ill-typed: it averages incommensurable burden states into a number that licenses nothing.


\subsection{7.2 Failure modes as relation failures (retained taxonomy)}
These are failures in the relation among occurrence, type, evidence, representation, and operation — not merely partition errors:


\begin{itemize}
\item \textbf{non-disclosure} — the system never notices that $m$ requires a distinction;

\item \textbf{orthemic underdetermination} — several consequential candidates remain open;

\item \textbf{misidentification} — $m$ is placed under the wrong ortheme;

\item \textbf{under-segmentation} — several orthemes in $O^*(m; A)$ are collapsed;

\item \textbf{over-segmentation} — distinctions are assigned that $m$ does not support;

\item \textbf{misplacement} — a valid description is attached to the wrong occurrence, component, level, or temporal version (the lineage failures of Section 2.5, including review-transport: a verdict formed about one artifact carried by the review process onto another);

\item \textbf{representational failure} — the correct ortheme is recognised but cannot be encoded;

\item \textbf{misrouting} — the correct ortheme is assigned but $m$ is sent to the wrong operation;

\item \textbf{validation failure} — the placement or repair is not genuinely tested;

\item \textbf{false closure} — $m$ is treated as fully resolved though consequential parts of its profile or its burden ledger remain open.

\end{itemize}
Each has a distinct remedy; they need not share a cause. Low confidence is not the only abstention trigger: broken provenance, an unavailable route, or a hard safety conflict can be decisive under a confident placement.


\subsection{7.3 Per-burden residual dispositions; false closure as type error}
\textbf{Definition 12 (Residual disposition).} Every burden of an episode carries, at closure time, exactly one disposition from


\begin{verbatim}
{ unresolved, deferred, transferred, owner-assigned, risk-accepted, validated-resolved }
\end{verbatim}
glossed: still open with no plan (unresolved); open with a declared later trigger (deferred); moved to another party with traceable ownership (transferred); waiting on a named owner's decision (owner-assigned); consciously left open under an accepted, recorded risk (risk-accepted); or closed on in-scope, current, sufficient evidence (validated-resolved). The six dispositions are \textbf{mutually exclusive by fiat} — each burden carries exactly one — with a precedence rule for the one genuinely overlapping pair: a burden that has been handed to a named decision-maker whose \emph{choice} is what is awaited is \textbf{owner-assigned}; a burden whose \emph{work} has been handed to another party is \textbf{transferred}. When both descriptions apply (work moved to a party who must also decide something), owner-assigned takes precedence, because the blocking condition is the decision. A disposition change is an event on the ledger, never a reinterpretation of the old record.

\textbf{Definition 13 (False closure, retained and typed).} False closure is a completion claim that collapses distinct dispositions — unresolved, deferred, transferred, owner-assigned, or risk-accepted — into "validated-resolved." In this revision it is a \textbf{type error in the record/schema sense} (no type-theoretic machinery is invoked): the completion claim quantifies over the burden ledger, and it is well-formed only when every burden's disposition admits it — a checkable well-formedness condition on the closure record, violated exactly when the claim and its own ledger disagree. "Done" uttered over a ledger containing a deferred burden is not an optimistic judgment; it is a claim whose type its own ledger refutes. Closure is legitimate exactly when every required burden is validated-resolved, or explicitly deferred/transferred/ owner-assigned/risk-accepted under an admissible policy, and the evidence supports the declared level of completion. Safe closure is reopenable by record; false closure accrues silent exposure.

\emph{Remark (demotion).} The prior draft's "orthemic debt" register is demoted to a speculative instrument: it remains unvalidated against ordinary defect counts, hazard registers, and technical-debt metrics, and no public evidence supports it. It is no longer part of the core proposal.


\subsection{7.4 The claim ledger and required success surfaces}
Closure quantifies over an explicit \textbf{claim ledger} $Q$: per claim, the record holds the proposition; the target occurrence and version $(\kappa, v)$ it is about; the property class of evidence that could support it (Section 4); the \textbf{required success surface}; the supporting evidence items by identifier; the warrant where relevant; the verification status; and explicit non-claims (what is \emph{not} being asserted, recorded to block later inflation).

The \textbf{required success surface} is the declared locus where the claim's success must be observable: a repair claim's success surface is the live behaviour of the current version, not the repository diff; a migration claim's surface is the serving system, not the migration log. A claim verified anywhere other than its declared surface is at best provisional. (The canonical synthetic instance: every check a repository can express passes while the user-visible surface is broken — the claim was verified off-surface.)


\subsection{7.5 Generalized ANDON: stop on risk, not on surprise (retained)}
A generalized ANDON event is a governed interruption when safe continuation or closure is not warranted under the current placement, evidence, model, route, or objective. Its signals are heterogeneous — observation novelty, intra-orthemma underdetermination, identity uncertainty, transition surprise, structural contradiction, evidence insufficiency or expiry, routing or capacity failure, validation failure, false-closure detection, cross-actor divergence (Section 2.6), unresolved metaorthemic conflict (Section 6.4), and safety conflict — and they are not the same signal. Stopping should track continuation risk, not anomaly:

\textbf{Proposition (ANDON sufficiency, retained framing).} If stopping or escalating is admissible with conditional risk at most $\theta_{\text{stop}}$, while every continuation action exceeds risk $\theta_{\text{stop}}$ or violates a hard constraint, then every risk-minimising admissible policy stops or escalates. \emph{(The stop/escalate action has lower conditional risk than every admissible continuation and meets the constraints, so no continuation is optimal.)}

The loop ends when residual obligations in the ledger have been reread — not when the original alarm disappears.


\subsection{7.6 Revision as a governed mode}
The prior draft's "revise" was one word. The record shows it needs its own machinery. A \textbf{revision operator}


\begin{verbatim}
ρ : (μ, lesson) ↦ μ′        (and analogously for repertoire revisions ρ : (O, lesson) ↦ O′)
\end{verbatim}
is governed by:


\begin{itemize}
\item \textbf{a formal trigger} — a recurring error class attributable to the current distinction or rule, not a single incident (one miss licenses a hypothesis, not a revision);

\item \textbf{a complexity penalty} — the revised rule or repertoire pays for its added distinctions against out-of-sample benefit (Definition 9's over-segmentation costs apply to rules as well as types);

\item \textbf{hysteresis} — revision and reversal thresholds differ, so the repertoire does not churn on boundary noise;

\item \textbf{impact-scoped reopening} — $\rho$ increments $\operatorname{ver}(\mu)$ and computes the set of \emph{affected} prior episodes: exactly those whose placement consulted the revised state-distinction or selector where the two editions disagree. Only those episodes' evidence status is downgraded from validated to provisional pending re-check. Blanket reopening of everything that ever ran under the old edition is rejected as unscoped;

\item \textbf{a second-order variant} — when a \emph{pattern across occurrences} (not any single case) shows the state space itself was under-specified, revision enters an enumeration mode: the alternatives are listed exhaustively against the record rather than patched one neighbor at a time. The synthetic worked example of Section 11.4 illustrates such a tranche: serial one-neighbor-per-review revision escalating to enumeration after repeated single-case patches fail to converge.

\end{itemize}

\bigskip\hrule\bigskip

\section{8. Episode Reification: Orthing Episodes and the Verdict Layer (formal addition 6)}

\subsection{8.1 The process type, the concrete token, and the reification embedding}
\textbf{Orthing} (gloss: the doing of apprehension-and-handling) is a process TYPE: the rule-governed, evidence-updating operation kind by which an orthemma is individuated, observed, assigned typed candidate families and an inferred profile, routed, validated, dispositioned, and revised under one or more metaorthemes. It stands to its runs as "compilation" stands to "the 14:32 build of commit \texttt{abc123}." An \textbf{orthing episode} $e$ is a concrete TOKEN: one dated, situated run, with an actor, a time, an input occurrence, and a result; two episodes are distinct occurrences even when inputs and outputs agree.

When a higher audit asks "was this placement made correctly, on adequate evidence, under an adequate rule?", the episode is itself apprehended, via an explicit \textbf{reification embedding} $\iota_n : E^{(n)} \hookrightarrow M^{(n+1)}$ — yesterday's classification act, re-cast as today's case. This is a modeling step, not set membership: an episode is a different kind of thing from a jar of powder and becomes a base occurrence of the higher audit only by being cast as one. The audit's orthemes at level $n+1$ are pathway state-types — evidence-sufficient placement, stale-rule placement, false closure — precisely the verdicts below. No separate noun ("meta-orthemma") is coined for the reified episode: against the paper's own anti-vacuity standard, the new noun changes no representation, evidence, routing, or prediction that the embedding does not already provide. The regress stops, as always, at the declared governance boundary.


\subsection{8.2 The episode signature (cited compactly)}
Following the core formalization, an orthing episode is a record


\begin{multline*}
e = \langle\, \mathrm{id};\ m, \kappa, v;\ x, H;\ \alpha, w, A, T, t;\\
\ \vec{\mu}, \mathrm{MetaTok}, \pi;\ \vec{C}, \hat{p};\ r;\\
\ \mathrm{estatus};\ \mathcal{Q};\ \delta;\ \mathrm{hand}_{\mathrm{in}},\\
 \mathrm{hand}_{\mathrm{out}};\ a, \mathrm{Succ} \,\rangle
\end{multline*}


— in words: which run this is; the concrete case with its identity key and version; what the episode saw, and the ordered typed evidence gathered (each item with property class, scope, provenance, validity); who or what executed ($\alpha$ — populated even for mechanical executors) and under what warrant ($w$, distinct from both evidence and executor identity); the declared analysis $A$ (identifier and version — the index against which result correctness is judged), its task $T = \operatorname{task}(A)$ retained as a separate readable component, and the time; the governing metaorthemic configurations with precedence; $\operatorname{MetaTok}(e)$ — the concrete \textbf{metaorthemmata} (Decision 0002): case-bound configuration tokens of those governing types, each recording or referencing its identity and lineage, type-and-version via $\operatorname{MetaInst}(\bar{\mu}, \mu)$, analysis compatibility $\operatorname{Compatible}(\bar{\mu}, A(e))$, occurrence anchor $(\kappa, v)$, governed component, case-specific binding map, scope with the claims that depend on it, policy/evidence-selector/instrument-and-calibration references, the binder with its binding warrant (kept distinct from the designated executor), binding time, and validity — referencing but never absorbing the episode's evidence and trace, and omitted entirely where no material case-specific binding exists (then V3c is inapplicable); the concrete policy executed under them; the typed candidate families and the inferred placement $\hat{p}$ (profile-valued, never coerced to a singleton ortheme); the route; the per-claim evidence-status map; the claim ledger with required success surfaces; the per-burden residual-disposition map; the incoming and outgoing handoff records; and the action with its labeled successor set. Components are defined \emph{where applicable}: a read-only classification episode has no action, no successors, an empty residual map — and the governance-derived required set $\operatorname{ReqPath}(e)$, together with a recorded \texttt{not-applicable} reason for every excluded verdict, keeps "deliberately none" distinct from "omitted." Where the pathway verdicts below must be adjudicated from the episode's own record, a bounded ordered trace of states and updates is also required, at a granularity governance declares — trivial cases carry no trace, and maximal logging is never a blanket requirement.

Three separations, stated once: the \textbf{episode is not its output} (episodes with identical outputs can differ in pathway and hence in every verdict); the \textbf{episode is not the policy} (the policy is repeatable and undated; reliability claims attach to the policy and its governing rules, correctness claims to the episode); a \textbf{placement inside an episode is not the episode} (the convenient one-placement notation $e \models_\mu (m : \hat{o})$ is derived, not definitional).


\subsection{8.2.1 Waking and somnic orthing}
A \textbf{waking orthing} takes up an occurrence under the analysis, evidence, governing configuration, claimant contracts, evaluator versions, and selector versions then in force. A later episode may reify that preserved orthing as its orthemma. \textbf{Meta-orthability} is the later episode's applicability and record-sufficiency gate; \textbf{somnic orthing} is the later placement of the prior orthing, pathway, residual, governing artifact, or relation. Thus: \textbf{waking orths experience; Somnus orths the available orthings of experience.} The earlier event history remains append-only: a later assessment may reference or supersede another assessment, but it cannot rewrite the target episode or insert later-discovered evidence into its historical evidence state.

The runtime-neutral contract distinguishes session, episode, occurrence, claim attempt, claimant-level orthability assessment, orthing, somnic assessment, and proposal identities; a conversational turn supplies none of these equalities. Privacy and source limits precede minimal occurrence capture, while versioned activation contracts gate claiming after capture. The record separately identifies evidence observed at episode time, used then, indexed but unused then, and discovered later. A somnic run selects newly unassessed or materially reopened anchors, queries an eligible historical corpus, and records the comparators actually used; a closed assessment is not recursively requeued without a material delta.


\begin{verbatim}
claim_status_ref: docs/decisions/0035-somnic-orthing-and-activation-contracts.md#somnus-claim-status
empirical_validation: not-established
performance: not-established
learning: not-established
generalization: not-established
internal_ontology: not-established
terminology_utility: not-established
terminology_adoption: not-established
runtime: not-implemented
collective_execution: not-implemented
\end{verbatim}
The bounded v0 operation is controlled residual-recurrence assessment. Fingerprint equality establishes only a structural recurrence candidate; support resolves through typed episode, session, orthing, actor, source-family, input-family, and occurrence-time records. A threshold requires distinct orthings and distinct episodes and is only a review trigger; the opportunity denominator is backed by explicit opportunity records, while independence uses a fixed declared policy. Suspected locus remains distinct from causal diagnosis or intervention. Assessment, intervention disposition, proposal, authorization, application, and later outcome evaluation remain separate. The schemas and fixtures under \texttt{schemas/\allowbreak{}} and \texttt{examples/\allowbreak{}somnus/\allowbreak{}} are offline conformance artifacts, not a live ledger, scheduler, analyzer, writeback engine, collective network, or learning result (Decision 0035).


\subsection{8.3 The verdict vector}
The verdicts diagnose distinct dimensions and are not identified with one another; they are NOT assumed pairwise logically independent — every definitional implication is declared explicitly, and the only one these definitions introduce is the claim-wise $\operatorname{V2b-T}_q \to \operatorname{V1}_q$ (full implication table: core formalization §4.1; Decision 0003, superseding the earlier "none entails another"). Verdict labels follow the normative registry (\texttt{docs/\allowbreak{}verdict-\allowbreak{}registry.\allowbreak{}yaml}, Decision 0004): semantic IDs — \texttt{RESULT\_CORRECT}, \texttt{EVIDENCE\_SUPPORT}, \texttt{PROCEDURE\_RELIABLE}, \texttt{TOKEN\_TRUTH\_LINKED}, \texttt{EVIDENCE\_CURRENT}, \texttt{GOV\_CONFIG\_ADEQUATE}, \texttt{GOV\_POLICY\_ADEQUATE}, \texttt{GOV\_TOKEN\_ADEQUATE}, \texttt{EXECUTION\_FAITHFUL}, \texttt{EX\_ANTE\_JUSTIFIED}, \texttt{ROUTE\_ADMISSIBLE}, \texttt{ROUTE\_QUALITY}, \texttt{CLOSURE\_TRUTHFUL}, \texttt{ROBUST\_NEIGHBORHOOD} — are authoritative in machine-readable records; the $\operatorname{V}\dots$ display aliases below are the prose forms. Pathway adequacy conjoins over the RESULT-FREE governed core $\operatorname{CorePath} = \{\operatorname{V2a}, \operatorname{V2b-P}, \operatorname{V2c}, \operatorname{V3a}, \operatorname{V3b}, \operatorname{V3c}, \operatorname{V3d}, \operatorname{V3e}, \operatorname{V4a}, \operatorname{V5}, \operatorname{V6}\}$ — V1 (result), V2b-T (factive), and route near-optimality are excluded. The required set $\operatorname{ReqPath}(e)$ is DERIVED from the declared analysis, episode shape, risk class, claims, and governance (never a discretionary list; "not tested" is never "not applicable"; every exclusion carries a recorded reason), and each verdict carries a status in \{pass, fail, undetermined, not-applicable\}: $\operatorname{PathwayAdequate}(e)$ iff every required verdict passes; $\operatorname{PathwayDefective}(e)$ iff one fails; $\operatorname{PathwayUndetermined}(e)$ when none fails but something required is unevaluated — a missing assessment is never silently a pass. (R3 honesty note: the repository ships the derivation as a machine-readable governance rule table with a per-verdict trace and an omission-attack fixture — \texttt{docs/\allowbreak{}governance-\allowbreak{}requirements.\allowbreak{}yaml}, \texttt{scripts/\allowbreak{}derive\_\allowbreak{}reqpath.\allowbreak{}py} — while $\operatorname{RequiredBy}$ \emph{in general} remains a governance-supplied parameterized interface, an acknowledged open parameter rather than a closed universal calculus; core formalization §4.1.)


% breakable-row-table: data-rows=13 total-rendered-characters=4535 max-row-rendered-characters=809
\par\medskip
\hrule height 0.8pt
\smallskip
% breakable-row: 1/13
\noindent\textbf{Verdict}: \textbf{V1 — result correctness}\par
\noindent\textbf{Question it answers}: Does the placed profile agree with $O^*(m; A(e))$ — the actual profile under the analysis recorded in the episode — at the governed level (or, weaker, is it route-sufficient with every placed claim true)? \textbf{Result-side: never a conjunct of pathway adequacy.}\par
\smallskip
\hrule height 0.4pt
\smallskip
% breakable-row: 2/13
\noindent\textbf{Verdict}: \textbf{V2a — evidential support}\par
\noindent\textbf{Question it answers}: Does the typed evidence, within its declared scopes, meet the declared standard for each placed claim?\par
\smallskip
\hrule height 0.4pt
\smallskip
% breakable-row: 3/13
\noindent\textbf{Verdict}: \textbf{V2b-P — configured-procedure truth-conduciveness} (pathway-side; NON-FACTIVE)\par
\noindent\textbf{Question it answers}: Does the procedure family actually instantiated here — under its governing configuration, applicable metaorthemmata, and execution mode — satisfy the predeclared reliability criterion over its DECLARED reference class (per-claim $\operatorname{RelSpec}_q$: reference class, stratum, metric, threshold, perturbation/comparison family, protocol, reliability evidence, version/validity)? The reference class and threshold are fixed independently of this episode's outcome; one current correct result cannot by itself establish it; \textbf{it does not entail V1} — a reliable configured procedure may produce a rare error. Default criterion: sensitivity; declared variants admissible.\par
\smallskip
\hrule height 0.4pt
\smallskip
% breakable-row: 4/13
\noindent\textbf{Verdict}: \textbf{V2b-T — token-level truth linkage} (result-side annotation; EXCLUDED from the pathway core)\par
\noindent\textbf{Question it answers}: Was THIS placed claim correct through the truth-relevant evidential mechanism rather than merely alongside it? FACTIVE and claim-wise: $\operatorname{V2b-T}_q \to \operatorname{V1}_q$; a profile-level reading requires every placed claim covered plus an explicit aggregation rule. Reportable for stopped-clock/Gettier diagnoses; excluded from $\operatorname{PathwayAdequate}$ precisely because its factivity would re-import result correctness. Non-factive token-local defects belong under V2a/V2b-P/V2c/V3a/V3b/V3c/V3d/V6, never under V2b-T alone.\par
\smallskip
\hrule height 0.4pt
\smallskip
% breakable-row: 5/13
\noindent\textbf{Verdict}: \textbf{V2c — evidence currentness}\par
\noindent\textbf{Question it answers}: Is each load-bearing evidence item current for $(\kappa, v)$ and of admissible provenance — not stale, not unsourced?\par
\smallskip
\hrule height 0.4pt
\smallskip
% breakable-row: 6/13
\noindent\textbf{Verdict}: \textbf{V3a — configuration adequacy}\par
\noindent\textbf{Question it answers}: Was each governing configuration adequate for the case's risk class — do its declared states separate what this case class can occupy, and can its selecting evidence actually discriminate them?\par
\smallskip
\hrule height 0.4pt
\smallskip
% breakable-row: 7/13
\noindent\textbf{Verdict}: \textbf{V3b — policy adequacy}\par
\noindent\textbf{Question it answers}: Was the meta-policy/procedure well-formed for the configuration it ran under? (A sound rulebook can carry an ill-formed procedure.)\par
\smallskip
\hrule height 0.4pt
\smallskip
% breakable-row: 8/13
\noindent\textbf{Verdict}: \textbf{V3c — governing-token adequacy} (Decision 0002)\par
\noindent\textbf{Question it answers}: Was every applicable concrete metaorthemma — the case-bound configuration token of a governing type (§8.2) — correctly instantiated, analysis-compatible, occurrence-anchored, correctly scoped to its claims, current, provenanced, and bound under authority? Per-token statuses are preserved alongside the episode-level conjunction; with no applicable token, V3c ∉ ReqPath(e) (zero-burden rule; status recorded not-applicable, with reason). Isolates: correct standard + sound policy + faithful execution + \textbf{defective case-specific binding} (wrong reference plane, wrong-role tolerance, expired calibration, wrong fixture or success surface). The Decision-0004 lettering matches the conceptual order: V3a → V3b → V3c (binding) → V3d (execution) → V3e.\par
\smallskip
\hrule height 0.4pt
\smallskip
% breakable-row: 9/13
\noindent\textbf{Verdict}: \textbf{V3d — executor fidelity}\par
\noindent\textbf{Question it answers}: Did the actor actually execute the procedure as written? (Adequate rule + adequate policy + infidelic executor routes to a different remedy: retrain or replace the executor, not rewrite the rule.)\par
\smallskip
\hrule height 0.4pt
\smallskip
% breakable-row: 10/13
\noindent\textbf{Verdict}: \textbf{V3e — ex-ante justification}\par
\noindent\textbf{Question it answers}: At decision time, on exactly the evidence then available, was the placement the reasonable one under the declared standard? (Blameless-at-the-time; indexed to decision time, unlike V2a/V2b-P/V2b-T, which are audit-time verdicts.)\par
\smallskip
\hrule height 0.4pt
\smallskip
% breakable-row: 11/13
\noindent\textbf{Verdict}: \textbf{V4a — route safety}\par
\noindent\textbf{Question it answers}: Was the route admissible under the constraints in force? Correct-route-unavailable is recorded as routing failure, not diagnostic uncertainty. (A finer near-optimality verdict may be recorded separately: safe ≠ best.)\par
\smallskip
\hrule height 0.4pt
\smallskip
% breakable-row: 12/13
\noindent\textbf{Verdict}: \textbf{V5 — closure truthfulness}\par
\noindent\textbf{Question it answers}: Does every residual carry an admissible disposition with traceable ownership, and does the completion claim match the ledger — no collapse of deferred/transferred/risk-accepted into "resolved"?\par
\smallskip
\hrule height 0.4pt
\smallskip
% breakable-row: 13/13
\noindent\textbf{Verdict}: \textbf{V6 — robustness}\par
\noindent\textbf{Question it answers}: Over the declared perturbation neighborhood of the episode — same policy, same rules, inputs from a declared perturbation family (version bumped, marker removed or spoofed, evidence reordered, near-identical sibling substituted) — is the V1 failure rate below tolerance? Paired metamorphic probes are V6's operational estimator.\par
\smallskip
\hrule height 0.8pt
\par\medskip

\subsection{8.4 The result × pathway matrix}
The matrix applies only after the pathway status resolves as adequate or defective; $\operatorname{PathwayUndetermined}$ episodes remain outside the four cells until audited. (Decision 0003: the earlier parenthetical committing token-level truth-connection to the conjunction is superseded — the incorrect-result/adequate-pathway cell is genuinely representable.)


% breakable-row-table: data-rows=2 total-rendered-characters=1697 max-row-rendered-characters=1115
\par\medskip
\hrule height 0.8pt
\smallskip
% breakable-row: 1/2
\noindent\textbf{Column 1}: \textbf{Result correct (V1)}\par
\noindent\textbf{PathwayAdequate}: \textbf{Nominal.} A release gate runs the behavioural suite against the exact commit it gates, on current artifacts, correctly bound (V3c where applicable), under a reliable configured procedure; it passes; the build is in fact good. Nothing to fix.\par
\noindent\textbf{PathwayDefective}: \textbf{Correct + defective (stopped clock / compensating error / governing-side luck).} The answer is right; the failure locus is V2b-P, V3a, V3b, V3c, V3d, or V6. The remedy targets the procedure, rule, or binding — not the verdict. Worked example below.\par
\smallskip
\hrule height 0.4pt
\smallskip
% breakable-row: 2/2
\noindent\textbf{Column 1}: \textbf{Result incorrect (¬V1)}\par
\noindent\textbf{PathwayAdequate}: \textbf{$\operatorname{AdequatePathError}(e) := \neg \operatorname{V1}(e) \wedge \operatorname{PathwayAdequate}(e)$} — the justified rare miss of a non-perfect but sufficiently reliable process used correctly. A triage router places a failing test under "infrastructure flake" because every available signal points there; the true cause is a rare race in new code: V3e holds, the procedure met its declared reference-class reliability (V2b-P), and the pathway IS certified adequate; \texttt{V2b-T} fails for the claim (it is false). Where V3e ∈ ReqPath(e), $\operatorname{JustifiedMiss}(e) := \operatorname{AdequatePathError}(e)$ (V3e's pass is already inside PathwayAdequate). Remedy: a new discriminating evidence source; no rule or executor blame. Orthemic adequacy does not establish moral, legal, institutional, or theological blamelessness.\par
\noindent\textbf{PathwayDefective}: \textbf{Compound failure.} A deploy gate reads a cached test result from the previous commit and approves the current one: wrong placement via a defective pathway (¬V1; V2c, V3a fail). Remedy: both the placement and the lineage machinery.\par
\smallskip
\hrule height 0.8pt
\par\medskip
Deterministic fixtures establishing all four resolved cells (plus the undetermined state, the stale-directive case, and safe-but-suboptimal routing) are given in the core formalization §4.2 (F1–F7) with machine-checkable encodings in \texttt{tests/\allowbreak{}verdict-\allowbreak{}fixtures.\allowbreak{}json}.


\subsection{8.5 Worked example: the stopped-clock validator}
A scan validator's rule is: \emph{pass iff the output log contains the marker string} \texttt{SCAN-CLEAN}. On today's artifact it passes, and the artifact happens to be genuinely compliant.


\begin{itemize}
\item \textbf{V1 holds:} the placed claim ("compliant") is true.

\item \textbf{V2a arguably holds by the letter:} the declared standard — marker presence — was met.

\item \textbf{V2b-P fails at the default sensitivity criterion:} over its declared reference class, had an artifact been non-compliant with the marker present (a tool version that always prints it; a spoofed log), the procedure as configured would still have supported the claim — the configured procedure is not truth-conducive. (\textbf{V2b-T also fails as a result-side diagnosis:} this claim was correct, but not \emph{through} the truth-relevant mechanism — the stopped-clock signature.)

\item \textbf{V6 fails:} the perturbation neighborhood exposes it directly — compliant-without-marker is rejected; non-compliant-with-marker is passed. The paired probes are exactly the metamorphic fixture pair.

\item \textbf{V3a fails:} the governing evidence configuration is inadequate for the risk class — its selecting evidence cannot discriminate the states it is supposed to separate.

\end{itemize}
This is the operational analogue of the stopped clock that shows the right time twice a day: correct result, defective pathway, high neighboring-case failure risk. The verdict layer expresses this on the episode record itself; a system tracking only result correctness has no object on which the defect is even statable, and will discover it only when a neighbor arrives.


\subsection{8.6 The novelty claim, stated narrowly}
The defensible claim is exactly this and no more: \textbf{the prior manuscript, and the operational discipline it describes, did not explicitly represent the concrete classification/handling episode as an auditable occurrence distinct from its result.} It is \emph{not} claimed that episodes are inexpressible elsewhere. Process reliabilism names precisely the correct-result-through-unreliable-process structure and the truth-connection idea — for \emph{beliefs}; it offers no operational record schema (no typed evidence, route, residual disposition, successor set, or perturbation estimator for agent and validator episodes). Provenance and audit-trail systems (W3C PROV, audit logs, ML lineage) are operational and token-level — who did what to which artifact when — but typically carry no candidate structure, plural inferred profile, per-claim evidence status, residual-disposition ledger, or verdict layer separating result correctness from pathway adequacy and robustness. Assurance practice (control-effectiveness vs outcome testing) and metamorphic testing each contain one verdict without the joint object. The residual contribution is the \textbf{integrated joint verdict layer on one auditable object}; whether that integration pays for itself over "audit the validator when suspicious" is an open empirical question (Section 13.2), not a settled fact.


\bigskip\hrule\bigskip

\section{9. Mechanical and Distributed Orthing}

\subsection{9.1 Consciousness is not required}
Nothing in the episode signature or the verdicts quantifies over awareness. Any rule-governed, evidence-updating executor — a predictive-text decoder resolving an ambiguous keypress sequence, a CI validator, a triage pipeline, a review chain, an institution — executes orthing, and its episodes bear all \emph{applicable} verdicts. Stopped-clock analogues arise mechanically: a decoder emits the right word from a frequency prior that would misfire on the neighboring input — correct message, defective pathway, no conscious subject anywhere. The actor field $\alpha$ is therefore populated with the mechanical executor — never left empty — and executor identity is distinct from any warrant or authority index. Semantic depth (hard-coded check < model-based inference over declared alternatives < reflective capacity to propose rule revisions) is a real axis of the \emph{executor}, orthogonal to every verdict on the \emph{episode}: shallow executors can run adequate pathways and deep ones defective pathways. Depth matters operationally in one place only: which revisions the executor may perform locally versus must escalate across the governance boundary.


\subsection{9.2 Distributed episodes: the token-level DAG}
One case handled across sensors, validators, agents, routers, a human sign-off, and downstream consumers is modeled as a finite \textbf{DAG} $\Gamma_E = (E, \rightsquigarrow)$ whose nodes are episodes (each with the full signature — its own actor, its own governing configuration) and whose edges are \textbf{typed}:


\begin{itemize}
\item \textbf{handoff} — a projection of one episode's output (a partial profile, an evidence item with its scope, a route, a disposition) appears in the input or evidence of the next; handoffs are the loci of transport error — a verdict crossing an edge without its scope and version is how stale evidence propagates;

\item \textbf{supersession} — a later episode revises an earlier one's conclusions;

\item \textbf{retry-of} — a later episode re-attempts an earlier one's task.

\end{itemize}
\textbf{Why a DAG despite retries and review loops:} nodes are \emph{dated tokens}; every edge is oriented from the earlier node to the later node with the semantics carried by its label (the later node of a supersession edge revises the earlier; the later node of a retry-of edge re-attempts the earlier's task), so the token graph is acyclic by construction. A rework loop does not add a back edge — it creates a \emph{new} episode node reached by a retry-of or supersession edge from its predecessor. The genuinely cyclic structure (the same policy re-entered; review returning work to an author role) lives at the TYPE/policy level, represented as repeated instantiation of the same policy, never as cycles among tokens.


\subsection{9.3 Composition conditions}
The graph composes into ONE boundary-level episode $e_\Gamma = \operatorname{comp}(\Gamma_E)$ iff:


\begin{enumerate}
\item \textbf{one case, one analysis} — all sub-episodes address $m$ or its successor-lineage under a single declared analysis $A$ (hence one task $T = \operatorname{task}(A)$; sub-episodes sharing a task but differing in tolerance, representation, or boundary do NOT compose without an explicitly declared fusion analysis);

\item \textbf{a declared boundary} — governance names the composite actor (pipeline, team, institution) and the composite governing configuration, including precedence across sub-episode rules;

\item \textbf{a declared fusion rule} — the composite placement, evidence-status map, and residual ledger are a stated aggregation of sub-outputs (e.g., profile union with per-claim status equal to the weakest status along its supporting path; merged ledgers with surviving ownership) — well-defined, not read off the last node.

\end{enumerate}
When condition 2 or 3 fails there is only the graph, and talk of "the system decided" is unwarranted composite attribution. Composite verdicts do not reduce to conjunctions: the composite can be correct while a sub-episode was wrong (a downstream validator caught it), and every sub-episode can be correct while the composite fails (each verdict true in scope, the composition transporting one out of scope). The composite level must be audited separately.


\bigskip\hrule\bigskip

\section{10. Multi-Actor and Competitive Settings}

\subsection{10.1 One occurrence, many evaluations}
The actor and analysis indices of Section 2.6 already imply that a single concrete occurrence supports several simultaneous evaluations: the same $(\kappa, v)$ is apprehended by different actors, each under its \textbf{own declared analysis} $A_\alpha$ — sharing frame components where the actors share a game or institution, and differing at least in task $T_\alpha = \operatorname{task}(A_\alpha)$ — each with its own inferred profile and candidate structure. This is a multi-analysis context: Definition 3's task-relative abbreviation is forbidden here, and ground truth is written in full, $O^*(m; A_\alpha)$, per evaluating analysis. This section develops the special case where the actors' interests diverge. In imperfect-information games the divergence begins one level earlier, at OBSERVATION: the one concrete deal $m$ is shared, but each actor observes only $x_\alpha = \Omega_\alpha(m)$ — the paper's own observation/occurrence distinction (§2.2), applied per actor. A player's information set is therefore an OBSERVATION, never part of the occurrence's identity key; perfect-information games are the special case of full, equal observation. Everything here is a \textbf{derived extension of the existing actor/analysis indices — not a seventh formal addition}; no new primitive is introduced.


\subsection{10.2 Target profiles}
\textbf{Definition 14 (Target profile set).} For actor $\alpha$ under its declared analysis $A_\alpha$ with task $T_\alpha = \operatorname{task}(A_\alpha)$, the target profile set $\mathcal{G}_{\alpha,A_\alpha} \subseteq \Pi_{A_\alpha}$ is the SET OF PROFILES that $\alpha$, under $T_\alpha$, aims to make some future occurrence in the lineage instantiate — each member of $\mathcal{G}_{\alpha,A_\alpha}$ is one complete profile (normative typing, aligned with the core formalization §5.3; a win-ortheme such as "White checkmates Black" labels the FAMILY of terminal profiles realizing it). $\mathcal{G}_{\alpha,A_\alpha}$ is the grounded instantiation at $\alpha$ of a PARAMETRIC ortheme schema $\operatorname{GoalSchema}(\alpha)$ ("win for α"): the schema is one form under role substitution; the grounded targets are distinct; and target-set overlap is a third, independent relation (empty for zero-sum win-sets, total for cooperation).

Three separations keep this well-typed:


\begin{itemize}
\item $\mathcal{G}_{\alpha,A_\alpha}$ is \textbf{not the current descriptive profile} $O^*(m; A_\alpha)$: the descriptive profile says what the occurrence \emph{is}; the target profile says what some successor should \emph{become}. A chess position's descriptive profile may include "material equality, White to move"; neither player's target is a description of the present board.

\item $\mathcal{G}_{\alpha,A_\alpha}$ is \textbf{not the objective/loss}: the loss is the task's evaluation function over outcomes; the target profile is the \emph{state-type content} of what the actor pursues — a set of orthemes over future occurrences, usable in placement and routing like any other type. (Objectives are excluded from metaorthemic governance, Section 6.2; target profiles, being typed state-sets under the actor/analysis index, are ordinary indexed objects.)

\item $\mathcal{G}_{\alpha,A_\alpha}$ is \textbf{not a metaortheme}: it names no governed component and declares no competing higher-order states about the machinery.

\end{itemize}

\subsection{10.3 Structural isomorphism is not identity of targets}
In symmetric games, a role substitution (swap the colors, relabel the players) maps one actor's target set onto the other's: the targets are \textbf{structurally isomorphic}. Declared prose abbreviation for this section and the next: the role labels \textbf{White} and \textbf{Black} name the formal actors $\alpha$ and $\beta$ respectively, with per-actor analyses $A_\alpha, A_\beta$ and tasks $T_\alpha = \operatorname{task}(A_\alpha)$, $T_\beta = \operatorname{task}(A_\beta)$. Isomorphism under role substitution is \emph{not} identity: $\mathcal{G}_{\alpha,A_\alpha} \neq \mathcal{G}_{\beta,A_\beta}$ even though each is the image of the other under the color swap. Conflating them is a placement error at the actor index — precisely the error a non-indexed formalism cannot even state, because it has only one profile slot per occurrence.


\subsection{10.4 Worked example: chess}
The shared concrete occurrence $m$ is \emph{this board position in this game}, where the occurrence's identity includes the position, the full move history, and the player to move (castling and en-passant rights are functions of that history). Both players apprehend the same $m$: one occurrence, two actor/analysis-indexed evaluations $\hat{p}_{A_\alpha,\alpha,t}(m)$ and $\hat{p}_{A_\beta,\beta,t}(m)$ (the §10.3 role-label convention: White = $\alpha$, Black = $\beta$), and two target profiles $\mathcal{G}_{\alpha,A_\alpha}$ (state-types realizing a win for White) and $\mathcal{G}_{\beta,A_\beta}$.

The example also calibrates what is and is not a metaortheme:


\begin{itemize}
\item the \textbf{transition rules} (legal moves) are part of the \emph{domain} — they define the successor structure of $M$;

\item the \textbf{utility} (win/draw/loss values) is the \emph{task};

\item \textbf{minimax search} is a \emph{policy} — both players may use similar policies over the shared transition rules;

\item \textbf{none of these is a metaortheme.} A metaorthemic configuration in a game would govern, for example, \emph{what counts as evidence of the opponent's information state} (in imperfect-information games: governed component E, competing states such as \{revealed, inferred-from-tempo, unknown\}, selecting evidence declared in advance), or \emph{when analysis may stop} (governed components D and V: a declared stopping standard for search — quiescence reached vs time-budget exhausted — whose variation changes which placements may be acted on). These pass the anti-vacuity conditions; the rules, the utility, and the search policy do not.

\end{itemize}

\subsection{10.5 Draws, cooperation, and conflict predicates}
\textbf{Draws exist:} not every terminal occurrence realizes a player's positive target. The two target sets do not partition the terminal states; there is a third region realizing neither. Any formalization that forces every terminal state into some actor's target set has mis-specified the game.

Over a set of actors with target profiles $\{\mathcal{G}_{\alpha_i,A_i}\}$, the predicates are well-typed across profile spaces — the target sets live in \emph{different} spaces ($\mathcal{G}_{\alpha,A_\alpha} \subseteq \Pi_{A_\alpha}$), so they are never compared by bare set intersection. With $\operatorname{Reach}(m)$ the occurrences reachable from $m$:


\begin{itemize}
\item \textbf{Compatibility:} $Compat_m(\mathcal{G}_1, \mathcal{G}_2)$ iff $\exists m' \in \operatorname{Reach}(m): O^*(m'; A_1) \in \mathcal{G}_1 \wedge O^*(m'; A_2) \in \mathcal{G}_2$ — one shared occurrence, evaluated under each analysis separately;

\item \textbf{Conflict:} $Conflict_m(\mathcal{G}_1, \mathcal{G}_2)$ iff no such reachable $m'$ exists — zero-sum chess targets conflict by construction;

\item \textbf{Cooperation (special case):} a shared target profile, $\mathcal{G}_{\alpha_1,A} = \mathcal{G}_{\alpha_2,A}$ under one shared analysis — the actors may still differ in evidence, roles, and routes (that is what the distributed-episode machinery of Section 9.2 is for), but their targets coincide.

\end{itemize}
Mixed-motive settings sit between: targets partially overlap, and the conflict/compatibility predicates apply region by region. All of this is bookkeeping over indexed profiles and lineage — the machinery of Sections 2 and 9 — which is why it is derived, not primitive.


\bigskip\hrule\bigskip

\section{11. Examples}
Each example is one architecture: orthemma → typed candidate families → typed evidence → validated placement → composed route → dispositioned residuals. The framework is substrate-neutral only as a design heuristic; it claims no shared mechanism across the domains. For each case the analyst must name the physical and social information carriers, boundary, decision rights, and validation.


\subsection{11.1 The build report}
The orthemma is \emph{this build report for this commit} — with identity key (the commit) and version (the artifact set actually built). The observation is the string "all tests pass." The candidate structure is typed: $C^{\mathrm{profile}}$ contains "behaviour genuinely exercised" and "files merely exist" as alternatives on the evidence-scope axis; $C^{\mathrm{id}}$ contains "report describes this commit" and "report describes the previous commit" (identity uncertainty, Section 3); $C^{\mathrm{warrant}}$ records whether the release is authorized, established, both, or neither. The discriminating evidence is typed: a behavioural channel (the suite demonstrably executed these paths on this artifact) versus a structural channel (the log line exists) versus a provenance channel (the log was produced by this run from this commit). A green check whose scope does not intersect the release claim is a mis-scoped pass and contributes nothing to closure. The route composes over resolved factors — release, repair, or reject — and closure quantifies over the burden ledger: a release with a deferred flake-investigation burden is closed \emph{with a deferred residual}, not "done."


\subsection{11.2 The Arabic word-token}
The orthemma is \emph{this word-token} on this page. Its candidate profile spans segmentation, lemma, lexeme, morphological pattern, clitics, syntactic role, and vocalisation — co-holding components, not alternatives: a future particle plus imperfect verb plus plural marking plus object clitic is one composed placement, and Arabic computational morphology explicitly separates tokenisation, lemmatisation, features, clitics, and contextual disambiguation. Root and lemma can be route-sufficient for dictionary lookup before every dependency is resolved — the factorized profile acts on resolved axes while the syntactic-role axis stays open with a declared evidence-to-resolve clause (more context). A fluent gloss does not validate a source-linked parse: the gloss is evidence on the semantic axis whose scope does not intersect the morphological claim.


\subsection{11.3 Clinical case (abstract only)}
A patient presentation is an orthemma whose factorized profile includes a symptom axis, a risk/urgency axis, a cause axis, and a severity axis. The symptom-axis placement can become route-sufficient — altering urgency and routing — while the cause axis stays open as an explicitly maintained candidate family with a discriminating test declared. Shared symptoms do not establish shared etiology (an evidence-grade configuration), and orthemic accuracy is not clinical or ethical goodness: a system can place correctly and route unjustly. The example is retained in this abstract form only; no clinical detail is load-bearing anywhere in the paper.


\subsection{11.4 Worked example: a launch pipeline (synthetic composite)}

\begin{quote}
\textbf{Synthetic composite — illustrative only; not evidence.}

This worked example is a synthetic composite assembled from the generic failure classes the paper's public schemas and fixtures already represent (identity, lineage-transport, warrant-binding, validator-scope, closure, and revision defects). It is patterned after the kind of engineering history that motivated the design, but it reports no auditable record, carries no observational weight, and is cited nowhere in this paper as support for any claim.

\textbf{Identity.} A storage slot is reused and the new occupant treated as the same file (identity key confused with location); elsewhere an ordinal position is mistaken for an identity key. Both are $C^{\mathrm{id}}$ failures — identity uncertainty misfiled as confidence about the wrong object. \textbf{Admission.} A case is admitted for handling on the basis of a stale checkout: evidence valid for $(\kappa, v)$ is applied to $(\kappa, v')$ — a lineage-transport violation, not a wrong judgment. \textbf{Authority-binding.} An authorization is present but not bound to the parameters actually used: the warrant's scope does not cover the executed action — a warrant-gate failure that no evidence channel could have expressed, because authorization is not evidence. \textbf{Settlement.} A validator returns green while its scope does not intersect the claim being settled (hash-valid read as semantically valid): a mis-scoped pass, detectable only where scopes are explicit. \textbf{Closure.} A terminal state is declared while residuals remain owner-gated: the completion claim collapses owner-assigned burdens into "resolved" — the false-closure type error, caught by the per-burden ledger. \textbf{Revision.} After repeated single-neighbor patches fail to converge, revision escalates to the enumeration mode of Section 7.6.

An operation without the union in advance would have to build each piece in response to each failure as this composite depicts. Whether that integration burden is real and common in actual practice is an empirical question reserved for Section 13 — this example illustrates the claim and cannot support it.

\end{quote}

\subsection{11.5 The stale steer (recovered directive without current force)}
A system recovers an earlier governing instruction — after context compaction, a crash, or a handover — and follows it. The recovery is textually faithful, the instruction was genuinely issued under proper authority, and it had been superseded in the interim. Five claims must be kept apart: that the directive \emph{exists}, that this copy is \emph{authentic}, that it was \emph{recovered} through a declared channel, that its issuer was \emph{authorized}, and that it is \emph{in force} now — and evidence for any one is not evidence for the next. The verdict layer needs no new machinery: currentness of the governing evidence fails (V2c), the bound governing token is not current (V3c), the executor is faithful (V3d passes — faithful execution under a defective governing token), and a closure claim of "acted under current instructions" fails V5. The full directive record, the five predicates, and the deterministic fixture (F6) are given in \texttt{examples/\allowbreak{}compaction-\allowbreak{}stale-\allowbreak{}steer.\allowbreak{}md} (Decision 0006).


\bigskip\hrule\bigskip

\section{12. Related Work}
Established constructs largely theorise the type side, and each facet of the present architecture has a nearest neighbor, conceded here facet by facet:


% breakable-row-table: data-rows=11 total-rendered-characters=1603 max-row-rendered-characters=269
\par\medskip
\hrule height 0.8pt
\smallskip
% breakable-row: 1/11
\noindent\textbf{Facet of this paper}: Plural, hierarchical profiles\par
\noindent\textbf{Nearest established neighbor}: multi-label and hierarchical classification (Tsoumakas \& Katakis 2007; Silla \& Freitas 2011)\par
\smallskip
\hrule height 0.4pt
\smallskip
% breakable-row: 2/11
\noindent\textbf{Facet of this paper}: Maintained candidate structure\par
\noindent\textbf{Nearest established neighbor}: POMDP belief states (Kaelbling, Littman \& Cassandra 1998); predictive-state representations (Littman, Sutton \& Singh 2001)\par
\smallskip
\hrule height 0.4pt
\smallskip
% breakable-row: 3/11
\noindent\textbf{Facet of this paper}: Abstention, escalation, quarantine\par
\noindent\textbf{Nearest established neighbor}: reject-option classification (Chow 1970); selective prediction (El-Yaniv \& Wiener 2010; Geifman \& El-Yaniv 2017); open-set recognition (Scheirer et al. 2013); OOD detection (Hendrycks \& Gimpel 2017)\par
\smallskip
\hrule height 0.4pt
\smallskip
% breakable-row: 4/11
\noindent\textbf{Facet of this paper}: Deciding when to stop investigating\par
\noindent\textbf{Nearest established neighbor}: value-of-information (Howard 1966) and sequential/optimal stopping (Wald 1947)\par
\smallskip
\hrule height 0.4pt
\smallskip
% breakable-row: 5/11
\noindent\textbf{Facet of this paper}: Residual obligations with owners\par
\noindent\textbf{Nearest established neighbor}: risk registers and risk-management guidelines (ISO 31000:2018); hazard logs\par
\smallskip
\hrule height 0.4pt
\smallskip
% breakable-row: 6/11
\noindent\textbf{Facet of this paper}: Placement statuses and transitions\par
\noindent\textbf{Nearest established neighbor}: ticket/workflow state machines (commodity practice; no single citable origin is load-bearing here)\par
\smallskip
\hrule height 0.4pt
\smallskip
% breakable-row: 7/11
\noindent\textbf{Facet of this paper}: Lineage, versions, audit records\par
\noindent\textbf{Nearest established neighbor}: W3C PROV-DM (Moreau \& Missier 2013); audit trails; ML lineage systems\par
\smallskip
\hrule height 0.4pt
\smallskip
% breakable-row: 8/11
\noindent\textbf{Facet of this paper}: Correct result through defective process\par
\noindent\textbf{Nearest established neighbor}: process reliabilism, for beliefs (Goldman 1979; the structure of Gettier 1963)\par
\smallskip
\hrule height 0.4pt
\smallskip
% breakable-row: 9/11
\noindent\textbf{Facet of this paper}: Perturbation probes for weak rules\par
\noindent\textbf{Nearest established neighbor}: metamorphic testing (Chen, Cheung \& Yiu 1998; Chen et al. 2018)\par
\smallskip
\hrule height 0.4pt
\smallskip
% breakable-row: 10/11
\noindent\textbf{Facet of this paper}: Rule adequacy vs outcome testing\par
\noindent\textbf{Nearest established neighbor}: assurance practice: tests of controls vs substantive procedures (ISA 330)\par
\smallskip
\hrule height 0.4pt
\smallskip
% breakable-row: 11/11
\noindent\textbf{Facet of this paper}: Type individuation and safe merger\par
\noindent\textbf{Nearest established neighbor}: sufficient statistics (Fisher 1922); causal states (Crutchfield \& Young 1989; Shalizi \& Crutchfield 2001); bisimulation and MDP state abstraction (Larsen \& Skou 1991; Givan, Dean \& Greig 2003); automata minimisation (Hopcroft 1971)\par
\smallskip
\hrule height 0.8pt
\par\medskip
The type-individuation mathematics is inherited outright. The proposed residual is exactly three things. (1) \textbf{The union:} no listed framework, in its standard individual form, supplies the others or names their union. (2) \textbf{The lifecycle:} the explicit, single-object connection of disclosure, evidence typing, representation, routing, validation, per-burden closure, and governed revision to the placement status of a concrete occurrence and its successors. (3) \textbf{The joint verdict layer:} process reliabilism names truth-connection for beliefs but has no operational record; provenance systems have operational records but no verdict layer; assurance and metamorphic testing each hold one verdict without the joint object. The orthing episode is one standardized object on which result correctness, evidential support, truth-connection, currentness, rule adequacy, policy adequacy, executor fidelity, ex-ante justification, route safety, closure truthfulness, and robustness are simultaneously definable. A bespoke composite could of course be extended with an equivalent record; the contribution claimed here is to state that integration once, generically, with a machine-checkable contract — an integration proposal and a standardization hypothesis, not an impossibility result. Whether the integration pays for itself is Section 13's question, not this section's claim.


\subsection{12.1 Sequential latent-variable inference (bounded related work)}
A neighbouring literature models cognition with sequential latent-variable machinery — hidden Markov models, clone-structured graphs, recurrent networks — in which ambiguous immediate observations are disambiguated by sequence-sensitive inference over model-internal latent states (George et al. 2021; Vasudeva Raju et al. 2024; Sun et al. 2025). One reported finding is structurally relevant here: in the cited study several compared model settings satisfy a declared similar-final-representational-structure criterion while only CSCG consistently matched the reported decorrelation order among the tested models under the reported evaluation, and high task performance is reported to occur without the global representation geometry that study measured. The bounded claim this paper draws is therefore about \emph{underdetermination}: \textbf{final performance and final representation do not by themselves determine the learning trajectory or the mechanism} — which is the same separation of result from pathway that Sections 8 and 13 make operational.

The boundary is fixed by Decision 0015 and is strict. A model's latent state $z$ is not an ortheme; its emitted observation is not the concrete occurrence; a posterior over latent states is not ground truth; and an internal representation's geometry (orthogonality, decorrelation, cluster separation) is neither necessary nor sufficient for orthemic distinctness or for pathway adequacy. Latent labels are identified only up to permutation and relabelling \textbf{absent declared semantic anchoring or alignment constraints}, so they carry no cross-model or cross-version meaning unless an explicit, validated alignment map is exhibited. The latent apparatus is an optional extension declared inside $A$: \textbf{the core claims of this framework remain statable without it}, though an extension may of course use its own declared vocabulary for extension-specific claims. This literature is cited as related work only. It validates nothing here, supplies no etymological support for the coined vocabulary, and bears on no theological question.


\subsection{12.2 A domain-specific application (bounded; not evidence, not part of the contribution claim)}
A separate, co-developed repository (\texttt{theislampill/\allowbreak{}daee-\allowbreak{}epistemics}, inspected read-only at a pinned commit) is a domain-specific noetic-diagnostic runtime whose objects can be \emph{typed through} this framework: an engagement is an immediate occurrence; public discourse is observation; an inferred noetic profile is a candidate placement, never the person's interior state; the governing distinctions are metaorthemes, their case bindings metaorthemmata, and the whole run an episode with a claim/residual ledger and a whole-state reread. That crosswalk (with a mandatory notation firewall and ten adversarial fixtures) is documented under \texttt{applications/\allowbreak{}daee-\allowbreak{}epistemics/\allowbreak{}}. Its role here is strictly bounded: it demonstrates that the architecture can \textbf{type a real designed system} (representability) and it \textbf{stress-tests} the framework's ability to keep state, observation, governing standard, application token, execution, closure, and restoration distinct. Because the two repositories share an owner and a related model lineage, any apparent convergence may reflect \textbf{co-development}; the application is therefore \textbf{not independent validation}, \textbf{not} evidence for any claim of this paper, and \textbf{not part of the contribution/novelty argument} — which remains an integration proposal adjudicated only by the unrun benchmarks of Section 13. It establishes nothing about fiṭrah, metaphysical orthability, a Necessary Being, or divine Speech.


\bigskip\hrule\bigskip

\section{13. Empirical Program}
Three evaluations, kept strictly separate. \textbf{Nothing in this paper is experimentally validated.} The current public evidence is exactly this: machine-checked internal agreement over the declared definitions, schemas, examples, and adversarial fixtures — which proves no mathematical consistency or completeness and no utility — and \textbf{no public observational dataset}. The first two programs below exist as frozen, deterministically validated, ready-to-run packets under \texttt{experiments/\allowbreak{}} (canonical per-packet state: \texttt{experiments/\allowbreak{}experiment-\allowbreak{}status.\allowbreak{}yaml}, Decision 0018); the terminology instrument awaits its owner-gated blind matching review. \textbf{None has run, and none is externally preregistered} — a Git freeze is not a preregistration, and registry submission is an owner/external act.


\subsection{13.1 Headline: the false-closure / selective-prediction benchmark (packet \texttt{FCSP-2} — preregistration-ready; not externally registered)}
Environments with planted inter-orthemma aliasing, planted identity uncertainty, and validator-scope traps. \textbf{Treatment:} an agent maintaining typed candidate families and per-burden closure (the machinery of Sections 5 and 7). \textbf{Baseline:} argmax placement with a confidence threshold. \textbf{Primary endpoints:} false-closure rate; the area under the risk–coverage curve (AURC) of the selective-prediction literature. (The metric name is corrected from an earlier draft's "risk–coverage AUROC", which named no established metric. Methodological criticisms of AURC — its dependence on the evaluation set's coverage grid and its incomparability across error rates — are addressed in the packet's frozen endpoint definitions by base-error-rate-normalized and per-family sensitivity analyses; any endpoint change is a new packet version, never a post-run edit.) Fixtures are laptop-scale and cheap; the benchmark targets the claimed residual directly (does the union behave differently from the composite?), not the vocabulary. The current frozen packet is \texttt{experiments/\allowbreak{}false-\allowbreak{}closure-\allowbreak{}selective-\allowbreak{}prediction-\allowbreak{}v2/\allowbreak{}} (packet \texttt{FCSP-2}; run harness with a mock adapter and public/hidden-key separation, neutral stimuli whose defects are inferable-not-stated, exact endpoints with paired permutation tests and Holm adjustment, mechanical decision-rule execution, a pre-run item-level sensitivity plan, deterministic smoke tests; status \texttt{READY\_TO\_RUN} / \texttt{PREREGISTRATION\_READY}; it passes the Decision 0022 methods gate). It supersedes the historical R6 packet \texttt{FCSP-1}, whose methods defects are recorded in Decision 0020. \textbf{No run has occurred; no result exists.}

\textbf{Decision rule — locally protocol-frozen in the packet (\texttt{DECISION-RULES.yaml}), four outcomes:}


\begin{itemize}
\item \textbf{supports incremental value} — both frozen primary contrasts favor the treatment at or beyond their minimum important effects, with no harm rule firing;

\item \textbf{does not yet support incremental value} — an adequately powered run in which neither the supports nor the harm rule fires (an adequately powered null IS this outcome);

\item \textbf{evidence of harm or failure} — accuracy degradation or structure-for-structure's-sake beyond the frozen harm margins;

\item \textbf{inconclusive} — only via a predeclared gate (failed-run rate, power floor, or an outcome-affecting recorded deviation), with the gate named.

\end{itemize}
An underpowered miss is never read as refutation; an unpowered tie is never read as adoption. This supersedes the prior draft's bare "no improvement over established labels" failure criterion, which conflated the second and third outcomes.


\subsection{13.2 The product fixture suite (separate product test)}
The current frozen packet is \texttt{experiments/\allowbreak{}episode-\allowbreak{}reification-\allowbreak{}v2/\allowbreak{}} (packet \texttt{ER-2}; five archetypes A1–A5 with four neutral surface variants each, both arms rendered label-free from one canonical fact list, a run harness, corrected scoring — completion by equality, semantic E5/traceability — complete paired inference, and mechanical decision execution; status \texttt{READY\_TO\_RUN} / \texttt{PREREGISTRATION\_READY}; it passes the Decision 0022 methods gate). It supersedes the historical R6 packet \texttt{ER-1}, whose scoring/leakage defects are recorded in Decision 0020. A fixture suite of the A1–A5 class exercises the product behaviour of an implementation: baseline-versus-treatment runs with pinned versions and run hygiene, culminating in \textbf{A5}, the metamorphic pathway fixture — a weak rule (pass iff the output contains a marker string) that returns the correct current result, paired with neighboring perturbations (correct-without-marker; incorrect-with-marker), scored to separate current-answer correctness from pathway adequacy. E5 carries no theory vocabulary and is the operational test of episode reification's practical delta over "audit the validator when suspicious" (Section 8.6). It is a product test; success or failure here transfers no verdict to the benchmark of 13.1 or the terminology test of 13.3.


\subsection{13.3 The terminology benchmark (explicitly secondary)}
The three-arm benchmark of the terminology evaluation spec: \textbf{Arm A} (ordinary language), \textbf{Arm B} (the operational distinctions expressed without neologisms), \textbf{Arm C} (the coordinated coined vocabulary with a matched primer), counterbalanced, exposure-matched, with construct-specific co-primary endpoints (false-closure detection; pathway-defect detection), predeclared minimum important effects, noninferiority and harm margins, and per-term three-outcome decision rules (frozen locally; externally registered only by an owner act). This benchmark is \textbf{secondary and cannot be rescued by 13.1}: a machinery win in the headline benchmark is an Arm-B-versus-Arm-A result about \emph{distinctions} and transfers nothing to the \emph{words}; only Arm C versus Arm B adjudicates the vocabulary. The converse holds too: a vocabulary win would not validate the machinery.


\subsection{13.4 Evidence status, stated plainly}
Machine-checked conformance: the definitions and separations of Sections 2–10 agree under the repository's declared schemas, examples, and adversarial fixtures — internal conformance over the declared cases, not a proof of consistency, completeness, or utility. Public observational: \textbf{none} — the private records that motivated the design are not independently auditable and carry no evidential weight in this paper (Section 17); the existence, recurrence, and base rates of the motivating failure patterns are publicly unestablished. Experimental: none. No stronger claim is made anywhere in this paper, and any sentence read as making one should be corrected against this section.


\bigskip\hrule\bigskip

\section{14. Terminology Status}
Every coined term in this paper — \textbf{orthemma, ortheme, metaorthemma, metaortheme, orthing,} and the companion-paper candidate \textbf{orthable} — is a \textbf{candidate} pending the matched benchmark of Section 13.3. None is adopted; none is retired. The paper's own governing principle applies:


\subsection{14.1 Formation and provenance boundary}
The vocabulary is \textbf{constructed but morphologically grounded}. Merriam-Webster supports established \texttt{orth-}, including its \textbf{correct or corrective} combining-form sense, and English \texttt{-eme}, whose linguistic-unit sense supplies an \textbf{analogy} rather than the project's operational definition. Smyth §§841.2 and 861.2 directly describe Greek \texttt{-ma} with stem \texttt{-mat-} as morphology for a \textbf{result or effect}; “object” or “instance” is only a bounded project gloss from the cited concrete examples. LSJ s.v. \texttt{ἐμός} supports Greek \texttt{ἐμά} as a deliberately superposed \textbf{possessive resonance}, \textbf{not a derivation} of \texttt{-ma}, \texttt{-emma}, or any project word.

The same spelling \texttt{ortheme} occurs in a prior, different \textbf{orthographic-unit sense} in Peter Constable's 17 June 1999 Unicode archive message and in Burridge and Blaxter (2020, p. 24). Those attestations do not supply this project's state-type sense. A bounded search that located no authoritative earlier hit for other project strings establishes no novelty, originality, priority, or first coinage. This provenance account \textbf{does not adopt} or retire any term and supplies no empirical utility result. Cross-domain usefulness remains a \textbf{benchmark-gated mnemonic or meta-schema hypothesis}, not universal primitives or a proven isomorphism.

\textbf{Design Principle 1 (Terminological utility, retained).} The vocabulary is warranted only if, under matched training time, attention, and access to evidence, it improves at least one predeclared outcome — placement accuracy, inter-rater agreement, defect discovery, repair traceability, transfer, calibration, or false-closure prevention — without unacceptable over-segmentation or governance cost, relative to established terminology. Adjudication follows the three-outcome rule: adopt on an adequately powered win; do not adopt yet on insufficient evidence; retire only on adequately powered equivalence or harm.

Two facts are recorded without tension. First: \textbf{every distinction in this paper is fully expressible in ordinary words} — occurrence and version, state-type, governing rule and the distinction it consults, handling episode and pathway verdict. The reader who strikes every coinage loses no content. Second: \textbf{paraphrasability does not settle the question} — an earlier claim that expressive redundancy of the words was "established analytically" conflated paraphrasability with redundancy and is retracted. "Gettier case" is the standing existence proof that a fully paraphrasable coined term can earn permanent keep through compression and coordination. Whether these terms do is exactly what the benchmark measures, and no argument in this paper substitutes for it. The distinctions are machine-checked for internal conformance and empirically untested; the vocabulary has no matched comparative evidence in either direction.


\bigskip\hrule\bigskip

\onecolumn
\appendix

\section{15. Limitations and Honest-Evidence Appendix}

\subsection{15.1 Limitations}
The decision-theoretic core is inherited, and the terminology may add load rather than value. Task specifications can be incomplete, contested, or imposed; exact equivalence is rarely estimable; approximate equivalence need not be transitive; local irreducibility depends on representation and merger families; system boundaries can be chosen opportunistically. Cross-domain analogies establish no shared mechanism. The formalism is conditional on objectives, losses, constraints, authority, and affected-party weights and cannot derive their legitimacy: a system can place every task-relevant ortheme correctly and still serve harmful ends — orthemic adequacy is not ethical goodness, and governance must permit contestation of the task, the loss, the authority, and the category itself. Orthemic debt is demoted to a speculative instrument (Section 7.3); the metaphysical and theological extensions are split to a companion paper and no result here depends on them. Metaorthemes remain higher-order control; the episode-verdict layer's practical delta is untested (fixture E5); and no study yet demonstrates incremental utility of anything in this paper. Acknowledged open formal parameters (R3): $\operatorname{RequiredBy}$ beyond the shipped machine-readable governance instance; evidence-class exhaustiveness (a hypothesis); fusion-mapping non-uniqueness; the $\Delta_A$ idealizations — the formal lane's honest ceiling is "internally conformance-checked specification over the declared definitions, schemas, examples, and adversarial tests — no proof of consistency, completeness, or utility," never mathematical completeness.


\subsection{15.2 Evidence tiers}

\begin{center}
\begin{tabular}{@{}p{0.313\linewidth}p{0.313\linewidth}p{0.313\linewidth}@{}}
\toprule
\textbf{Tier} & \textbf{Content} & \textbf{Status} \\
\midrule
\textbf{Conceptual/definitional} & The proposed distinctions and typed relationships of Sections 2–10; the split metaortheme normal form; verdict separability; the DAG/composition conditions; the derived multi-actor extensions & Machine-checked internal agreement over the declared definitions, schemas, examples, and adversarial fixtures (deterministic); \textbf{not} a proof of mathematical consistency or completeness; establishes no empirical claim \\
\textbf{Public observational} & — & \textbf{None currently supplied.} The private records that motivated the design are non-evidential design provenance (Section 17) \\
\textbf{Experimental} & Machinery benefit (13.1); episode-reification delta (13.2); vocabulary utility (13.3) & \textbf{None.} All three designed, none run \\
\bottomrule
\end{tabular}
\end{center}
Three honesty notes close the paper. \emph{No public observational dataset exists:} the design-history materials are private and not independently auditable, and they carry no evidential weight here — the existence, recurrence, and base rates of the motivating failure patterns are all publicly unestablished. \emph{Deterministic fixtures are not proofs:} they demonstrate conformance over the declared cases only. \emph{Nothing has been independently replicated:} no external party has yet examined any tier beyond this public repository itself.


\bigskip\hrule\bigskip

\section{16. Conclusion}
This paper proposed the orthemma–ortheme system as an integration discipline for the token side of classification and handling work. Its center of gravity is one relation — a concrete occurrence instantiating repeatable operational state-types relative to a declared, versioned analysis — and one further object built on that relation: the orthing episode, an auditable record on which result correctness and pathway adequacy are separately adjudicable. Around these, six formal additions close gaps the prior text could only gesture at: versioned identity with labeled lineage makes "right finding, wrong copy" a statable error; analysis/actor/time indexing makes ground truth explicit and disagreement diagnosable; typed, scoped, expiring evidence makes the green-but-mis-scoped validator visible; factorized candidates with route composition make partial action well-formed; per-burden closure makes false closure a checkable record defect; and episode reification with the registry-normalized verdict vector makes the stopped clock, the justified rare miss, the defective binding, and the stale directive all first-class, machine-checkable cases (fixtures F1–F7).

What the paper deliberately does not claim is as much a result as what it claims. Every facet has a conceded neighbor; the type-individuation mathematics is inherited; a bespoke composite could represent the architecture, so the contribution is the standardized integration, not an impossibility claim; the private records that motivated the design are not independently auditable, carry no evidential weight, and validate nothing; the coined vocabulary is fully paraphrasable and strictly benchmark-gated; and the three designed studies — the false-closure benchmark (packet FCSP-2) and the episode-reification test (packet ER-2), both locally frozen and ready to run, and the terminology comparison (instrument awaiting its blind matching review) — have not been run, and none is externally preregistered. The framework's standing is therefore: internally conformance-checked (deterministic fixtures — not a consistency proof), design-motivated, empirically unvalidated. The next honest step is not more theory but the owner-authorized execution of Section 13's frozen packets — externally registered at run time — and the framework has been built so that a negative result would be recordable, on its own terms, as exactly what it is.


\section{17. Data and Materials Availability}
All public materials of this project — the formal core, this manuscript, the multi-actor note, machine-readable schemas (hardened, with negative and mutation test suites) and worked examples, deterministic fixtures and validators (including the governance-derivation, cross-record-semantics, Qurʾān-locus, type/token, and terminology-matching validators added in R3), the verdict and notation registries with their generated surfaces, decision records, the R3 sourcing ledgers, and the designed (unrun) terminology-evaluation packets — are available in the public repository (\texttt{theislampill/orthemology} on GitHub), with a SHA-256 release manifest under \texttt{docs/\allowbreak{}provenance/\allowbreak{}} and byte-reproducible draft PDFs under \texttt{artifacts/\allowbreak{}}. \textbf{Not available:} the private internal design-history records (a cross-project defect collection and a single-project longitudinal engineering record) that motivated the framework. They contain private working material, are not rights-cleared for publication, and are therefore not independently auditable; they are non-evidential design provenance — \textbf{no claim in this paper rests on them} (Sections 11.4, 13.4, 15.2). \textbf{No public observational dataset accompanies this paper.} No human-subject data exists; no experiment has been conducted.


\twocolumn

\section{References}

\begin{itemize}
\item Chen, T. Y., Cheung, S. C., \& Yiu, S. M. (1998). \emph{Metamorphic testing: a new approach for generating next test cases.} Technical Report HKUST-CS98-01, Department of Computer Science, HKUST.

\item Chen, T. Y., Kuo, F.-C., Liu, H., Poon, P.-L., Towey, D., Tse, T. H., \& Zhou, Z. Q. (2018). Metamorphic testing: a review of challenges and opportunities. \emph{ACM Computing Surveys}, 51(1), 4:1–4:27. doi:10.1145/3143561

\item Chow, C. K. (1970). On optimum recognition error and reject tradeoff. \emph{IEEE Transactions on Information Theory}, 16(1), 41–46. doi:10.1109/TIT.1970.1054406

\item Crutchfield, J. P., \& Young, K. (1989). Inferring statistical complexity. \emph{Physical Review Letters}, 63(2), 105–108. doi:10.1103/PhysRevLett.63.105

\item El-Yaniv, R., \& Wiener, Y. (2010). On the foundations of noise-free selective classification. \emph{Journal of Machine Learning Research}, 11, 1605–1641.

\item Fisher, R. A. (1922). On the mathematical foundations of theoretical statistics. \emph{Philosophical Transactions of the Royal Society A}, 222, 309–368. doi:10.1098/rsta.1922.0009

\item Geifman, Y., \& El-Yaniv, R. (2017). Selective classification for deep neural networks. \emph{NeurIPS 30}.

\item Geifman, Y., Uziel, G., \& El-Yaniv, R. (2019). Bias-reduced uncertainty estimation for deep neural classifiers. \emph{ICLR}. (Derivation of AURC/E-AURC.)

\item Gettier, E. L. (1963). Is justified true belief knowledge? \emph{Analysis}, 23(6), 121–123. doi:10.1093/analys/23.6.121

\item Givan, R., Dean, T., \& Greig, M. (2003). Equivalence notions and model minimization in Markov decision processes. \emph{Artificial Intelligence}, 147(1–2), 163–223. doi:10.1016/S0004-3702(02)00376-4

\item Goldman, A. I. (1979). What is justified belief? In G. S. Pappas (Ed.), \emph{Justification and Knowledge} (pp. 1–23). D. Reidel. doi:10.1007/978-94-009-9493-5\_1

\item Hendrycks, D., \& Gimpel, K. (2017). A baseline for detecting misclassified and out-of-distribution examples in neural networks. \emph{ICLR}.

\item Hopcroft, J. (1971). An n log n algorithm for minimizing states in a finite automaton. In \emph{Theory of Machines and Computations} (pp. 189–196). Academic Press.

\item Howard, R. A. (1966). Information value theory. \emph{IEEE Transactions on Systems Science and Cybernetics}, 2(1), 22–26. doi:10.1109/TSSC.1966.300074

\item International Auditing and Assurance Standards Board (2009). \emph{ISA 330: The Auditor's Responses to Assessed Risks.} IFAC.

\item International Organization for Standardization (2018). \emph{ISO 31000:2018 Risk management — Guidelines.}

\item Kaelbling, L. P., Littman, M. L., \& Cassandra, A. R. (1998). Planning and acting in partially observable stochastic domains. \emph{Artificial Intelligence}, 101(1–2), 99–134. doi:10.1016/S0004-3702(98)00023-X

\item Larsen, K. G., \& Skou, A. (1991). Bisimulation through probabilistic testing. \emph{Information and Computation}, 94(1), 1–28. doi:10.1016/0890-5401(91)90030-6

\item Littman, M. L., Sutton, R. S., \& Singh, S. (2001). Predictive representations of state. \emph{NIPS 14}.

\item Moreau, L., \& Missier, P. (Eds.) (2013). \emph{PROV-DM: The PROV Data Model.} W3C Recommendation, 30 April 2013. https://www.w3.org/TR/prov-dm/

\item Merriam-Webster. “Orth-.” \emph{Merriam-Webster.com Dictionary.} https://www.merriam-webster.com/dictionary/orth- (accessed 26 July 2026).

\item Merriam-Webster. “-eme.” \emph{Merriam-Webster.com Dictionary.} https://www.merriam-webster.com/dictionary/-eme (accessed 26 July 2026).

\item Smyth, H. W. (1920). \emph{A Greek Grammar for Colleges.} American Book Company, §§841.2, 861.2.

\item Liddell, H. G., Scott, R., \& Jones, H. S. \emph{A Greek-English Lexicon}, s.v. \texttt{ἐμός}.

\item Constable, P. (1999, June 17). “Re: Amerindian Characters.” \emph{Unicode Mail List Archive.}

\item Burridge, J., \& Blaxter, T. (2020). Using spatial patterns of English folk speech to infer the universality class of linguistic copying. \emph{Physical Review Research}, 2, 043053, p. 24. doi:10.1103/PhysRevResearch.2.043053

\item Scheirer, W. J., Rocha, A., Sapkota, A., \& Boult, T. E. (2013). Toward open set recognition. \emph{IEEE TPAMI}, 35(7), 1757–1772. doi:10.1109/TPAMI.2012.256

\item Shalizi, C. R., \& Crutchfield, J. P. (2001). Computational mechanics: pattern and prediction, structure and simplicity. \emph{Journal of Statistical Physics}, 104, 817–879. doi:10.1023/A:1010388907793

\item Silla, C. N., \& Freitas, A. A. (2011). A survey of hierarchical classification across different application domains. \emph{Data Mining and Knowledge Discovery}, 22, 31–72. doi:10.1007/s10618-010-0175-9

\item Traub, J., Bungert, T. J., Lüth, C. T., Baumgartner, M., Maier-Hein, K. H., Maier-Hein, L., \& Jaeger, P. F. (2024). Overcoming common flaws in the evaluation of selective classification systems. \emph{NeurIPS 37}. arXiv:2407.01032

\item Tsoumakas, G., \& Katakis, I. (2007). Multi-label classification: an overview. \emph{International Journal of Data Warehousing and Mining}, 3(3), 1–13. doi:10.4018/jdwm.2007070101

\item Wald, A. (1947). \emph{Sequential Analysis.} John Wiley \& Sons.

\item Wetzel, L. (2018). Types and Tokens. In E. N. Zalta (Ed.), \emph{The Stanford Encyclopedia of Philosophy.} https://plato.stanford.edu/entries/types-tokens/

\end{itemize}
Machine-readable database with per-claim verification status: \texttt{references/\allowbreak{}orthemology.\allowbreak{}bib}; current sourcing state starts at \texttt{docs/\allowbreak{}sourcing/\allowbreak{}CURRENT-\allowbreak{}SOURCING-\allowbreak{}LEDGER.\allowbreak{}md} (Decision 0019).


\bigskip\hrule\bigskip

\section{Glossary of Proposed Terms (all candidates; none adopted)}

\begin{itemize}
\item \textbf{Orthemma.} A concrete situated occurrence considered as something to be apprehended, carrying an identity key and version. The token pole.

\item \textbf{Ortheme.} A repeatable operational state-type an orthemma instantiates, relative to a declared, versioned analysis. The type pole. Task-relative wording is shorthand only after one analysis with that task has been fixed.

\item \textbf{Instantiation relation $\operatorname{Inst}_A$.} $(m, o) \in \operatorname{Inst}_A$: $m$ instantiates $o$ under the declared, versionable analysis $A$ (with $T = \operatorname{task}(A)$); the fibre $O^*(m; A)$ is the actual profile. $\operatorname{Inst}_T$ / $O^*_T(m)$ are licensed local shorthand once a single $A$ with $\operatorname{task}(A) = T$ has been fixed; forbidden in multi-analysis contexts (Definition 3).

\item \textbf{Typed candidate families $C^{\mathrm{id}}, C^{\mathrm{profile}}, C^{\mathrm{cause}}, C^{\mathrm{route}}, C^{\mathrm{warrant}}$.} Open alternatives per uncertainty axis, with exclusivity marking and an evidence-to-resolve clause.

\item \textbf{Placement / apprehension.} Assigning an inferred profile / the whole encounter-to-disposition process.

\item \textbf{Route-sufficient apprehension.} Enough of the profile to route safely while other axes stay open; ≠ identity-complete resolution.

\item \textbf{Labeled successor set $\operatorname{Succ}_a(m)$.} The zero, one, or many occurrences an action creates; placement validity does not transport across an edge without re-evidence.

\item \textbf{Metaorthemma (candidate term; object adopted by Decision 0002).} The episode-local, case-bound configuration TOKEN of a metaortheme: $\operatorname{MetaInst}(\bar{\mu}, \mu)$; binds case-specific values (reference frame, tolerance value, instrument + calibration, fixture, success surface, scope) within the declared analysis ($\operatorname{Compatible}(\bar{\mu}, A(e))$); its binder-with-warrant is distinct from the executor; judged by V3c; omitted where no material binding exists. The WORD remains a benchmark-gated candidate (ordinary-language equivalent: "instantiated governing-configuration token"); the OBJECT is part of the episode record.

\item \textbf{Metaortheme.} A metaorthemic configuration $\langle g, S_\mu, \operatorname{select}_\mu, \operatorname{prov}(\mu), \operatorname{ver}(\mu)\rangle$ paired with a separable meta-policy; governed component in \{repertoire, individuation, evidence, disclosure, routing, validation, warrant-classification\}; objectives, the task, and the governance meta-level excluded.

\item \textbf{Residual disposition.} Per burden: unresolved / deferred / transferred / owner-assigned / risk-accepted / validated-resolved.

\item \textbf{False closure.} A completion claim collapsing other dispositions into "validated-resolved"; a type error against the burden ledger.

\item \textbf{Orthing / orthing episode.} The process type of rule-governed, evidence-updating apprehension-and-handling / one dated, situated, record-valued run of it, bearing the verdict vector.

\item \textbf{Verdict vector (V1–V6, registry-normalized).} Result: V1 correctness, V2b-T factive truth linkage. Pathway core: V2a support; V2b-P procedure reliability; V2c currentness; V3a configuration, V3b policy, V3c token-binding, V3d execution, V3e decision-time adequacy; V4a route safety; V5 closure truthfulness; V6 robustness. Advisory: V4b route quality. Semantic IDs and aliases: \texttt{docs/\allowbreak{}verdict-\allowbreak{}registry.\allowbreak{}yaml} (Decision 0004).

\item \textbf{Generalized ANDON event.} A governed interruption when continuation or closure is not warranted under the current placement, evidence, route, warrant, or ledger.

\item \textbf{Target profile $\mathcal{G}_{\alpha,A_\alpha}$.} The state-types an actor aims to make some successor occurrence instantiate, under the actor's own analysis $A_\alpha$; distinct from the descriptive profile, from the objective, and from any metaortheme.

\end{itemize}
\emph{End of revised draft. Prior version: the original manuscript (see \texttt{docs/\allowbreak{}provenance/\allowbreak{}document-\allowbreak{}history.\allowbreak{}md}). Companion formal text: \texttt{theory/\allowbreak{}orthemic-\allowbreak{}core-\allowbreak{}formalization.\allowbreak{}md}. Terminology evaluation protocol: \texttt{terminology/\allowbreak{}}.}

\emph{R7B bounded extensions (candidate; require fresh review before merge).} Two application-level extensions deepen the framework without changing the core: a dynamic-orthing / representation-learning account (four update levels, world-vs-learner edges, ablation-based ortheme admission) in \texttt{applications/\allowbreak{}latent-\allowbreak{}state-\allowbreak{}orthing/\allowbreak{}} (Decision 0024), and a meta-noetic memetics + sound-corrective-dynamics account (the represented-standard bearer, the ecology, and the G1 route-ranked descent adjudication) in \texttt{applications/\allowbreak{}daee-\allowbreak{}epistemics/\allowbreak{}} (Decision 0025). The six-layer firewall — formal core, dynamic extension, noetic application, school-neutral metaphysics, Atharī identification, empirical programs — is mapped in \texttt{docs/\allowbreak{}architecture/\allowbreak{}ORTHEMOLOGY-\allowbreak{}LAYER-\allowbreak{}MAP.\allowbreak{}md}. None of it is empirical validation or a theological proof.*


\twocolumn
\bibliographystyle{plainnat}
\bibliography{../../../references/orthemology}
\end{document}
